\documentclass[journal]{IEEEtran}
\usepackage{amsmath,amssymb,mathtools}
\usepackage{amsthm}
\usepackage{booktabs,array}
\usepackage{algorithm}
\usepackage{algpseudocode}
\usepackage{enumitem}
\usepackage{microtype}
\usepackage[hidelinks]{hyperref}
\newtheoremstyle{plainn}{}{}{\itshape}{}{\bfseries}{.}{.5em}%
  {\thmname{#1}\thmnumber{ #2}\thmnote{ (#3)}}
\theoremstyle{plainn}
\newtheorem{thm}{Theorem}
\newtheorem{lem}[thm]{Lemma}
\newtheorem{prop}[thm]{Proposition}
\newtheorem{cor}[thm]{Corollary}
\theoremstyle{remark}
\newtheorem{remark}{Remark}
\renewcommand{\qedsymbol}{QED}

\begin{document}
\title{Deciding functional observer feasibility: compression, certificates, and
parameterized algorithms}
\author{}
\maketitle

\begin{abstract}
We study the parameterized complexity of deciding whether a functional observer of prescribed
order exists for a linear system with rational data and a prescribed block of observer rows. The two
feasibility conditions are coupled: the observer matrix and its product with the system matrix are
built from the same free rows. Compressing to the common kernel of the output matrix and of its
product with the system matrix reduces the free rows to a subspace of a compressed row-pair space, in
which the two conditions become a row-space inclusion and the constant rank of a pencil. On this
reduction we build a finite algebraic certificate for the spectral condition, compressed by
Cauchy--Binet to a Gram matrix whose size is governed by $\ell=\operatorname{rank}(AN)$ rather than by
the number of maximal minors, so no maximal minor is enumerated. Two bounds follow: $\mathcal L^{O(q^2)}$
bit operations for every fixed order $q$, and
$\exp\bigl(O(\ell^2\log(\ell+2))\bigr)\operatorname{poly}(\mathcal L)$ bit operations, fixed-parameter
in $\ell$, which does not involve the number $r$ of free rows, unlike the coordinate count
$\kappa_{\rm obs}=r\operatorname{rank}[N\ AN]$ that a chart parameterization suggests. The case of no
free rows is decided exactly in polynomial time, and the decision problem and its finite threshold
version lie in the existential theory of the reals. These are parameterized upper bounds; the
unrestricted complexity remains open.
\end{abstract}

\begin{IEEEkeywords}
functional observer; parameterized complexity; real algebraic geometry; polynomial matrix; fixed-parameter tractability; decision problems.
\end{IEEEkeywords}

\section{Introduction}\label{sec:intro}

Estimating a linear function of the state, rather than the whole state, is often what a control loop
needs, and the order of a functional observer can be far below that of a Luenberger observer for the
same system. The design question is quantitative: given the system and the functionals to be
estimated, what is the smallest number of observer rows for which a functional observer exists?

For a linear system with measured output, \cite{fer10} reduced that question to two rank conditions on
the observer matrix: with the observer rows written as a block $L$ containing the prescribed rows
$L_0$, one condition is algebraic and one is spectral. The characterization has since been completed:
\cite{dar25} gives a unified algebraic framework, and \cite{fer26a,fer26b} determine, order by order,
which orders are attainable --- \cite{fer26a} through functional observability indices, \cite{fer26b}
through a nullspace representation of the algebraic condition and a matrix-pencil form of the spectral
one, together with a procedure that parameterizes the nullspace over Grassmann charts with
$(n-p-q)(q-r)$ parameters per chart and tests the conditions chart by chart.

What these works do not provide is a cost. Neither states a running-time or a bit-complexity bound, and
the chart-by-chart procedure is not by itself an algorithm with a guarantee: both
$(n-p-q)(q-r)$ and the number of charts grow with $n$ even when the prescribed order $q$ and the number
$r$ of free rows are held fixed. Enumerating charts therefore establishes neither polynomial time at
fixed order nor a quantity that could control the cost.

This paper studies that complexity. Three things are new. First, the conditions are made finite: the
free rows act only through a subspace of a compressed row-pair space, and the spectral condition is
replaced by a finite algebraic certificate of bounded degree, so that no maximal minor is enumerated
and the fixed-order algorithm needs no chart enumeration; the fixed-parameter algorithm does enumerate
charts, but at most $4^{\ell}$ of them. Second, the parameter that controls the cost is identified,
namely $\ell=\operatorname{rank}(AN)$, an invariant of the system and of the prescribed rows alone that
does not involve $r$, whereas the parameter suggested by the chart parameterization, $r\,d_0$, does.
Third, the cost is bounded in the bit model: polynomial in the input length at fixed order, and
fixed-parameter in $\ell$.

Adjacent design problems are hard, which makes the question natural but does not answer it: selecting
sensors that render a system functionally observable is NP-hard \cite{zha25}, and static output feedback
with a spectral predicate is NP-hard \cite{lof26,bar26} and, in its real feasibility form,
$\exists\mathbb R$-complete \cite{liu26}. None of these results transfers: in sensor placement the
output matrix is itself a design variable, and the predicates concern a spectral abscissa rather than a
constant rank that has to hold for every complex value of the spectral variable.

The approach is to make the decision problem finite-dimensional. Restricting the observer rows to the
common kernel of the output matrix and of its product with the system matrix removes exactly the
directions those two matrices already account for, and turns the two conditions into a row-space
inclusion and into the constant rank of a pencil. The free rows then act only through the subspace $W$
they span in the compressed row-pair space. That reading gives two quantities: the dimension
$\kappa_{\rm obs}=rd_0$ of the coordinates of the free rows, and the dimension
$\ell=\operatorname{rank}(AN)\le k$ of the space in which the feasible $W$ live, which does not involve
$r$.

\paragraph*{Contributions and organization}\label{sec:contrib}
The results form a chain, and the paper follows it: each item states what is established and where.

\begin{enumerate}
\item \emph{The reduction} (Section~\ref{sec:compressed}). The two conditions are compressed to
$O(q^2)$ variables independently of the state dimension, and the free rows are reduced to a subspace
$W$ of a compressed row-pair space, in which they become a row-space inclusion and the constant rank of
a pencil.
\item \emph{The parameter} (Section~\ref{sec:dimension}). The cost is governed by
$\ell=\operatorname{rank}(AN)$, which does not involve $r$. We do not claim that $\ell$ is the only
parameter that could control the cost; we identify it as one that is independent of $r$.
\item \emph{The certificate} (Section~\ref{sec:certificate}). The spectral condition becomes a finite
algebraic certificate, compressed by Cauchy--Binet to a matrix of size controlled by $\ell$; no
maximal minor is enumerated.
\item \emph{The fixed-order algorithm} (Section~\ref{sec:fixedq}). $\mathrm{DEC}(q)$ is decided in
$\mathcal L^{O(q^2)}$ bit operations for every fixed order $q$, without enumerating charts.
\item \emph{The fixed-parameter algorithm} (Section~\ref{sec:fpt}). $\mathrm{DEC}$ is decided in
$\exp\bigl(O(\ell^2\log(\ell+2))\bigr)\operatorname{poly}(\mathcal L)$ bit operations. The exact
polynomial-time algorithm for the case of no free rows (Section~\ref{sec:zero}) is a separate special
case, not a consequence of this bound.
\item \emph{Classification} (Section~\ref{sec:existsr}). The prescribed-order decision problem and its
finite threshold version lie in the existential theory of the reals.
\end{enumerate}

\paragraph*{Roadmap}
Both complexity theorems are reached in the same three stages, and the intermediate results are named
here once, so that the later sections can also be read out of order. (i) \emph{Reduction}:
Lemmas~\ref{lem:compression}--\ref{lem:attainable} compress the two conditions to the row space of
$[H\ H_A]$, and Lemmas~\ref{lem:sandwich}--\ref{lem:W-conditions} make them depend on a subspace $W$
alone, which is what makes the search finite. (ii) \emph{Structure}: Lemma~\ref{lem:struct} confines a
feasible $W$ to $\omega+1\le\dim W\le2\omega$, Section~\ref{sec:US} bounds $\omega$ by
$\ell=\dim S$, and Lemma~\ref{lem:monotone} gives
monotonicity. (iii) \emph{Certificate}: Lemma~\ref{lem:bezout} and Lemma~\ref{lem:window} turn the
spectral condition into the Gram equation $Z^{\mathsf T}Zx=e_0$ of Proposition~\ref{prop:gram}, which is
the object both algorithms test. Theorem~\ref{thm:fixedq} then counts the cost at fixed order and
Theorem~\ref{thm:fpt} the cost in $\ell$; the only technical lemmas outside this chain are
Lemmas~\ref{lem:rationalbasis} and~\ref{lem:coeff}, which bound the coefficient sizes.

Section~\ref{sec:conclusions} collects what is not claimed and the evidence behind each statement;
Appendix~\ref{app:bits} records the bit-size estimates and Appendix~\ref{app:encoding} the
existential-real encoding. The optimization and construction problems, and a lower bound on the
minimum order, are treated in a companion paper.

\section{Preliminaries}\label{sec:prelim}

\subsection{Notation and conventions}\label{sec:notation}

We take $C$ to have full row rank, $p=\operatorname{rank}C$: otherwise $C$ is replaced by a basis of
its row space, which leaves $\operatorname{row}[C;CA]$ unchanged and hence leaves \eqref{eq:5a} and
\eqref{eq:5b} unchanged, since both are rank identities on matrices whose rows are augmented by $C$ and
$CA$. The integer $p$ in the order range \eqref{eq:range} is this rank.

Consider the linear time-invariant system
\begin{equation*}
\dot{x}=Ax+Bu,\quad y=Cx,\quad z=L_0x,
\end{equation*}
where $A\in\mathbb Q^{n\times n}$, $C\in\mathbb Q^{p\times n}$ and
$L_0\in\mathbb Q^{r_0\times n}$.

Table~\ref{tab:notation} collects the notation; the rows of the observer matrix are named there.
Let $q\ge r_0$ be the prescribed observer order, with $q\le n$, and set $r=q-r_0$. The
observer rows are the rows of the block
\begin{equation*}
L=\begin{bmatrix}L_0\\R\end{bmatrix},\qquad R\in\mathbb R^{r\times n},
\end{equation*}
where $L_0$ is fixed while $R$ collects the $r$ free rows. Throughout we call the $r_0$ rows of
$L_0$ the \emph{prescribed rows}, the $r$ rows of $R$ the \emph{free rows}, and the $q$ rows of $L$
the \emph{observer rows}; only real observers are considered.

All given rational data are encoded in binary, and $\mathcal L$ denotes the total input
bit length. The free parameter $R$ is not part of the given input data and is therefore
not counted in $\mathcal L$. State vectors are column vectors, and matrices act on them by
the usual left multiplication.

Let $P(\lambda)$ be a rational matrix pencil of generic rank $\omega$, and call its
$\omega\times\omega$ minors maximal minors. Ranks at spectral points $\lambda\in\mathbb C$
are always computed over $\mathbb C$. Because the maximal minors have rational
coefficients, their gcd is a nonzero constant in $\mathbb C[\lambda]$ exactly when it is a
nonzero constant in $\mathbb Q[\lambda]$; this equivalence is used without further
mention.

\subsection{The observer conditions}\label{sec:conditions}

A functional observer of order $q$ is a matrix $L=\begin{bmatrix}L_0 \\ R\end{bmatrix}\in\mathbb R^{q\times n}$ such
that
\begin{equation}\label{eq:5a}
\operatorname{rank}\begin{bmatrix}LA\\ L\\ CA\\ C\end{bmatrix}
=\operatorname{rank}\begin{bmatrix}L\\ CA\\ C\end{bmatrix},
\end{equation}
\begin{equation}\label{eq:5b}
\operatorname{rank}\begin{bmatrix}\lambda L-LA\\ CA\\ C\end{bmatrix}
=\operatorname{rank}\begin{bmatrix}L\\ CA\\ C\end{bmatrix},
\quad \forall\lambda\in\mathbb C.
\end{equation}
These are the Fernando--Trinh--Jennings conditions, numbered (5a) and (5b) in
\cite{fer10}. Condition~\eqref{eq:5a} is algebraic; \eqref{eq:5b} is spectral, and it is
required at every $\lambda\in\mathbb C$, in particular at all finite eigenvalues of $A$.

\subsection{Decision, optimization and construction}\label{sec:problems}

Three problems are posed for \eqref{eq:5a} and \eqref{eq:5b}, in increasing order of
demand.

\emph{Decision.} $\mathrm{DEC}$ takes $(A,C,L_0)$ and $q$ and asks whether some $R$
satisfies \eqref{eq:5a} and \eqref{eq:5b}; no witness is required. This is the problem
studied in this paper.

\emph{Optimization.} $\mathrm{OPT}$ asks for the minimum feasible order
\begin{equation*}
q_{\min}=\min\{q:\mathrm{DEC}(q)\},
\end{equation*}
where $\mathrm{DEC}(q)$ abbreviates $\mathrm{DEC}(A,C,L_0,q)$.

\emph{Construction.} $\mathrm{FIND}$ asks for a witness $R$ at the minimum feasible
order, $q=q_{\min}$.

The threshold problem is
\begin{equation}\label{eq:mindec}
\mathrm{MINDEC}(Q)=\bigvee_{q=r_0}^{\min(Q,n)}\mathrm{DEC}(q),
\end{equation}
a finite disjunction of decision instances. The algorithms below decide $\mathrm{DEC}$ and their
affirmative certificates contain a witness. Only $\mathrm{DEC}$ is studied here: $\mathrm{OPT}$ and
$\mathrm{FIND}$, together with the monotonicity and saturation of $\mathrm{DEC}$ in $q$ that make them
tractable, are treated in the companion paper \cite{companion2}.

\begin{table}[!t]
\centering
\caption{Notation. The first block is fixed by the input, the second by the compression, the third by
the subspace reduction of Section~\ref{sec:subspace}.}
\label{tab:notation}
\footnotesize
\begin{tabular}{@{}ll@{}}
\toprule
Symbol & Meaning \\
\midrule
$K_0=\ker[C;CA]$ & common kernel of $C$ and $CA$ \\
$k=\dim K_0$ & its dimension \\
$N$ & basis matrix of $K_0$ \\
$M=[N\ AN]$ & compressed row-pair matrix \\
\midrule
$H=LN$, $H_A=LAN$ & compressed observer blocks \\
$\omega=\operatorname{rank}H$ & compressed order \\
$P(\lambda)=\lambda H-H_A$ & compressed pencil \\
\midrule
$\mathcal T=\operatorname{row}M$ & attainable row-pair space \\
$d_0=\dim\mathcal T$ & its dimension \\
$\kappa_{\rm obs}=r\,d_0$ & number of free-row coordinates \\
$S=\operatorname{row}(AN)\subseteq\mathbb R^k$ & space of the second blocks \\
$\ell=\operatorname{rank}(AN)\le k$ & its dimension; the FPT parameter \\
$W=\operatorname{row}[H\ H_A]$ & row space of the compressed pairs \\
$W_0=\operatorname{row}(L_0M)$ & contribution of the prescribed rows \\
$s_0=\dim W_0$, $m=d_0-s_0$ & dimensions of $W_0$ and of $\mathcal T/W_0$ \\
$\mathcal V=\mathcal T\cap(S\times S)$ & admissible ambient space \\
$m'=\dim\mathcal V-s_0$ & chart dimension \\
$\pi_1,\pi_2,\pi_\lambda$ & projections of $\mathbb R^{2k}$ \\
\bottomrule
\end{tabular}
\end{table}

Two published bounds confine the orders that have to be examined. Let $q_0$ be the algebraic
lower bound of \cite{fer26a,fer26b}: the smallest order at which \eqref{eq:5a} can be satisfied,
computed from the functional observability indices by a greedy rank recursion, hence in polynomial
time. Every feasible order satisfies
\begin{equation}\label{eq:range}
q_0\le q\le n_0-p,\qquad n_0=\operatorname{rank}\bigl[C;CA;\dots;CA^{n-1}\bigr],
\end{equation}
see \cite{fer26a,fer26b}; for an observable pair $n_0=n$. The upper end is not merely a bound in the
row-slot model: by Lemma~\ref{lem:monotone} the predicate is monotone in $q$, so every order at or
above $n_0-p$ is feasible as soon as that one is, and no information is lost by stopping there. In
terms of the quotient dimension the range \eqref{eq:range} reads
$q_0-r_0\le t\le n_0-p-r_0$, and Lemma~\ref{lem:struct} adds the independent bound
$t=\dim(W/W_0)\le\dim W\le2\ell$.

\section{The reduction}\label{sec:compressed}\label{sec:dimension}

Both observer conditions depend on the free rows only through the products $LN$ and $LAN$, where $N$ is a basis matrix of the common kernel of $C$ and $CA$. This section performs that compression and rewrites \eqref{eq:5a} and \eqref{eq:5b} in terms of the compressed data.

\subsection{The common kernel and the compression}\label{sec:compress}

Let
\begin{equation*}
K_0=\ker\begin{bmatrix}C\\ CA\end{bmatrix},\qquad k=\dim K_0,
\end{equation*}
and let $N\in\mathbb R^{n\times k}$ be a basis matrix for $K_0$. Define
\begin{equation*}
H=LN\in\mathbb R^{q\times k},\quad H_A=LAN\in\mathbb R^{q\times k},\quad \omega=\operatorname{rank}H .
\end{equation*}
The compressions below rest on a single rank identity.

\begin{lem}[Rank identity]\label{lem:rankid}
Let $P\in\mathbb C^{m\times n}$ and let $N_P\in\mathbb C^{n\times r}$ be a basis matrix for $\ker P$. Then, for every $M\in\mathbb C^{p\times n}$,
\begin{equation}\label{eq:rankid}
\operatorname{rank}\begin{bmatrix}P\\ M\end{bmatrix}
=\operatorname{rank}P+\operatorname{rank}(MN_P).
\end{equation}
\end{lem}

\begin{proof}
Let $f(x)=xN_P$. Since $\operatorname{col}N_P=\ker P$, the kernel of $f$ is $\operatorname{row}P$,
and $\operatorname{row}(MN_P)$ is the image of the restriction of $f$ to $\operatorname{row}M$, so
rank--nullity gives $\operatorname{rank}[P;M]=\operatorname{rank}P+\operatorname{rank}(MN_P)$; this
is the standard rank identity for a stacked matrix \cite[Ch.~0]{hj13}.
\end{proof}

\subsection{The compression and the pencil criterion}\label{sec:compression}

\begin{lem}[Compression and pencil criterion]\label{lem:compression}
With $H$, $H_A$ and $\omega$ as above, the following statements hold.

(i) Condition \eqref{eq:5a} holds if and only if $\operatorname{row}H_A\subseteq\operatorname{row}H$.

(ii) If \eqref{eq:5a} holds, then there exists $F\in\mathbb R^{q\times q}$ such that $H_A=FH$.

(iii) Condition \eqref{eq:5b} holds if and only if
\begin{equation*}
\operatorname{rank}(\lambda H-H_A)=\omega,\qquad \forall\lambda\in\mathbb C .
\end{equation*}
This part uses only the rank identity. If in addition \eqref{eq:5a} holds, then
$P(\lambda)=(\lambda I-F)H$ has rank at most $\omega$ everywhere, and the condition is equivalent to
the $\omega\times\omega$ maximal minors of $P(\lambda)$ having no common complex root, equivalently to
their greatest common divisor in $\mathbb C[\lambda]$ being a nonzero constant.

(iv) If \eqref{eq:5a} and \eqref{eq:5b} hold, then $\operatorname{rank}H_A=\operatorname{rank}H$.
\end{lem}

\begin{proof}
Set $P_0=[C;CA]$. Each part is one rank computation on the stacked matrices of \eqref{eq:5a} and
\eqref{eq:5b}: with $P=P_0$ and $N$ a basis of $\ker P_0$, Lemma~\ref{lem:rankid} replaces every
stacked rank by $\operatorname{rank}P_0$ plus the rank of a compressed block, and what remains is to
read off the resulting row inclusion and to propagate it to $\lambda H-H_A$.

\emph{(i)} The rows of $P_0$ occur among the rows of both matrices below, so Lemma~\ref{lem:rankid} with $P=P_0$, $M=\begin{bmatrix}LA\\L\end{bmatrix}$ and $M=L$ gives
\begin{equation*}
\begin{aligned}
\operatorname{rank}\begin{bmatrix}LA\\ L\\ CA\\ C\end{bmatrix}
&=\operatorname{rank}P_0+\operatorname{rank}\begin{bmatrix}H_A\\ H\end{bmatrix},\\
\operatorname{rank}\begin{bmatrix}L\\ CA\\ C\end{bmatrix}
&=\operatorname{rank}P_0+\operatorname{rank}H .
\end{aligned}
\end{equation*}
Up to row permutations the left-hand sides are the two sides of \eqref{eq:5a}, so \eqref{eq:5a}
holds exactly when $\operatorname{rank}[H_A;H]=\operatorname{rank}H$, that is, when
$\operatorname{row}H_A\subseteq\operatorname{row}H$.

\emph{(ii)} By this row inclusion every row of $H_A$ is a linear combination of the rows of $H$, so
$H_A=FH$ for the matrix $F$ whose $i$-th row is the coefficient vector of the $i$-th row of $H_A$.

\emph{(iii)} With $M(\lambda)=[\lambda L-LA;CA;C]$ and $J=[L;CA;C]$, and with $N$ as above,
$M(\lambda)N=[P(\lambda);0;0]$ and $JN=[H;0;0]$. Lemma~\ref{lem:rankid} with $P=P_0$,
$M=\lambda L-LA$ and $M=L$ gives
$\operatorname{rank}M(\lambda)=\operatorname{rank}P_0+\operatorname{rank}P(\lambda)$ and
$\operatorname{rank}J=\operatorname{rank}P_0+\operatorname{rank}H$, hence \eqref{eq:5b} is equivalent
to
\begin{equation}\label{eq:5b-1}
\operatorname{rank}P(\lambda)=\operatorname{rank}H=\omega,\ \forall \lambda\in\mathbb C.
\end{equation}
Assume now \eqref{eq:5a}. By $H_A=FH$, $P(\lambda)=\lambda H-H_A=(\lambda I-F)H$, so
$\operatorname{row}P(\lambda)\subseteq\operatorname{row}H$ and therefore
$\operatorname{rank}P(\lambda)\le\omega$ for all $\lambda\in\mathbb C$. Combining \eqref{eq:5b-1} with
this bound, condition \eqref{eq:5b} becomes the requirement that some $\omega\times\omega$ minor of
$P(\lambda)$ be nonzero at every $\lambda$; a common complex root of the maximal minors is therefore
exactly a value of $\lambda$ at which all of them vanish, which gives the two equivalent forms in the
statement, and for $\omega=0$ the condition is vacuous.

\emph{(iv)} Assume \eqref{eq:5a} and \eqref{eq:5b}. By (iii), \eqref{eq:5b-1} holds for every
$\lambda$; at $\lambda=0$, where $P(0)=-H_A$,
$\operatorname{rank}H_A=\operatorname{rank}P(0)=\operatorname{rank}H$. \qedhere
\end{proof}

\subsection{The space of attainable row pairs}\label{sec:free}

Each free row $r_i\in\mathbb R^{1\times n}$ of $R$ determines the row pair
\begin{equation*}
(H_i,(H_A)_i)=r_i\,[N\ \ AN] \in\mathbb R^{1\times 2k}.
\end{equation*}
Consequently, this attainable pair is uniquely specified by its coordinates relative to a basis of the subspace of attainable pairs
\begin{equation*}
\begin{gathered}
\mathcal T=\operatorname{row}[N\ \ AN]\subseteq\mathbb R^{1\times 2k},\\
d_0=\dim\mathcal T=\operatorname{rank}[N\ \ AN].
\end{gathered}
\end{equation*}
Therefore, a rational basis of $\mathcal T$ allows each attainable pair generated by a free row to be parameterized by exactly $d_0$ real coordinates.

\begin{lem}[Attainable row pairs and the effective dimension]\label{lem:attainable}
Every free row pair lies in $\mathcal T$, and writing the $r$ free rows in a basis of $\mathcal T$ represents them by
\begin{equation}\label{eq:kappa}
\kappa_{\rm obs}=r d_0
\end{equation}
coordinates. Moreover, when $r\ge1$,
\begin{equation}\label{lem3-2}
k\le d_0\le rd_0=\kappa_{\rm obs}.
\end{equation}
\end{lem}

\begin{proof}
Since $N$ has full column rank,
\begin{equation}\label{lem3-1}
k=\operatorname{rank}N\le\operatorname{rank}[N\ AN]=d_0 ,
\end{equation}
and $d_0\le rd_0$ for $r\ge1$, so \eqref{lem3-2} holds. The count $\kappa_{\rm obs}=rd_0$ is a
coordinate count, not the dimension of an image, since many $R$ give the same row pairs. The case
$r=0$ is excluded here and handled in Section~\ref{sec:zero}.
\end{proof}

\begin{remark}[Basis invariance]
If $N\mapsto NQ$ with nonsingular $Q\in\mathbb R^{k\times k}$, then
$[N\ \ AN]\mapsto[N\ \ AN]\operatorname{diag}(Q,Q)$ with $\operatorname{diag}(Q,Q)$ nonsingular, so
$\mathcal T$, $d_0$ and $\kappa_{\rm obs}=rd_0$ are unchanged; the free-row coordinates undergo the
nonsingular change $r_i\mapsto r_i\operatorname{diag}(Q^{-1},Q^{-1})$.
\end{remark}

\subsection{The subspace reduction}\label{sec:subspace}

The free rows enter the two conditions only through the row space of the compressed
pair. Writing the row pair of a free row as a row of $M=[N\ AN]$, let
\begin{equation*}
W=\operatorname{row}[H\ H_A]\subseteq\mathbb R^{2k},
\end{equation*}
so that $W_0=\operatorname{row}(L_0M)\subseteq W\subseteq\mathcal T=\operatorname{row}M$.
Let
\begin{equation}\label{eq:ell}
\begin{gathered}
d=\dim\mathcal T,\qquad s_0=\dim W_0,\qquad m=d-s_0,\\
\ell=\operatorname{rank}(AN)\le k,\qquad S=\operatorname{row}(AN)\subseteq\mathbb R^k ,
\end{gathered}
\end{equation}
and let $\pi_1(h,g)=h$, $\pi_2(h,g)=g$, $\pi_\lambda(h,g)=\lambda h-g$ be the projections of
$\mathbb R^{2k}=\mathbb R^k\times\mathbb R^k$. The parameter $\ell$ is computed in polynomial time: $N$ comes from exact rational elimination on
$[C;CA]$ and $\ell$ is the rank of the rational matrix $AN$. Because \eqref{eq:5b} quantifies over complex
$\lambda$, the projection $\pi_\lambda$ is applied to the complexification
$W_{\mathbb C}=W\otimes_{\mathbb R}\mathbb C\subseteq\mathbb C^{2k}$, and dimensions are over
$\mathbb C$; a real matrix has the same rank over $\mathbb R$ and over $\mathbb C$.

\begin{lem}[Sandwich and realizability]\label{lem:sandwich}
For $L=[L_0;R]$ one has $W_0\subseteq W\subseteq\mathcal T$ and
$\dim(W/W_0)\le r$. Conversely, every subspace $W$ with these three properties is realized by a
free block $R$ with exactly $t=\dim(W/W_0)$ rows.
\end{lem}

\begin{proof}
$W=\operatorname{row}(L_0M)+\operatorname{row}(RM)$ gives $W_0 \subseteq W\subseteq\mathcal T$ and $\dim(W/W_0)\le\operatorname{rank}(RM)\le r$. 
For the converse, the map $u\mapsto uM$ is onto $\mathcal T$; choose preimages $u_1,\dots,u_t$ of a basis of
$W/W_0$ and let $R$ consist of those $t$ rows followed by $q-r_0-t$ zero rows. Then
$W_0+\operatorname{row}(RM)=W$.
\end{proof}

\begin{lem}[The conditions depend only on $W$]\label{lem:W-conditions}
Condition \eqref{eq:5a} holds if and only if $\pi_2(W)\subseteq\pi_1(W)$, and condition
\eqref{eq:5b} holds if and only if $\dim\pi_\lambda(W)=\dim\pi_1(W)$ for every
$\lambda\in\mathbb C$.
\end{lem}

\begin{proof}
By Lemma~\ref{lem:compression}(i), \eqref{eq:5a} holds exactly when
$\operatorname{row}H_A\subseteq\operatorname{row}H$, which is $\pi_2(W)\subseteq\pi_1(W)$ since
$W=\operatorname{row}[H\ H_A]$. By Lemma~\ref{lem:compression}(iii) and $\operatorname{rank}H=\omega$,
\eqref{eq:5b} holds exactly when $\operatorname{rank}(\lambda H-H_A)=\omega$ for every $\lambda$, which
is $\dim\pi_\lambda(W)=\dim\pi_1(W)$, since $\pi_1(W)=\operatorname{row}H$ has dimension
$\omega=\operatorname{rank}H$.
\end{proof}

\begin{lem}[Structural bound]\label{lem:struct}
Let $W$ be feasible and $\omega=\dim\pi_1(W)>0$. Then
\begin{equation}\label{eq:struct}
\omega+1\le\dim W\le2\omega .
\end{equation}
\end{lem}

\begin{proof}
For the upper bound, Lemma~\ref{lem:W-conditions} with $\lambda=0$ gives
$\dim\pi_0(W)=\dim\pi_1(W)=\omega$; since $\pi_0(W)=\pi_2(W)$ up to sign,
$U:=\pi_2(W)=\pi_1(W)$ and $W\subseteq U\times U$, so $\dim W\le2\dim U=2\omega$. Also
$\dim W\ge\dim\pi_1(W)=\omega$, which combines with the lower bound below.

For the lower bound, if $\dim W=\omega$ then $\pi_1|_W$ and $\pi_2|_W$ are isomorphisms and $W$ is the
graph of a linear $F:U\to U$, that is $g=Fh$ on $W$; the same conclusion follows from
$\dim\pi_2(W)=\dim\pi_0(W)=\omega$, so this step does not depend on the previous paragraph. Over
$\mathbb C$ the complexification $F_{\mathbb C}$ has an eigenvalue $\lambda_0$, and
$\pi_{\lambda_0}(W)=(\lambda_0I-F)U$ has dimension $<\omega$, contradicting
Lemma~\ref{lem:W-conditions}. The hypothesis $\omega>0$ is used exactly here, to have an eigenvalue;
at $\omega=0$ feasibility forces $H=H_A=0$ and $\dim W=0$, so the bound does not survive.
\end{proof}

\paragraph*{The projection lies in $S$}\label{sec:US}
For a feasible $W$, Lemma~\ref{lem:compression}(iv) gives
$\operatorname{rank}H_A=\operatorname{rank}H$, while Lemma~\ref{lem:compression}(i) gives the row
inclusion $\operatorname{row}H_A\subseteq\operatorname{row}H$; the two spaces therefore coincide, and
\begin{equation*}
\pi_1(W)=\pi_2(W)=\operatorname{row}H_A\subseteq\operatorname{row}(AN)=S .
\end{equation*}
This is the inclusion that lets every feasible object be written in the $\ell$ coordinates of
$S$: it is used for the Cauchy--Binet compression of Section~\ref{sec:cb} and for the chart
coordinates of Section~\ref{sec:chartcoords}.

\begin{lem}[Monotonicity]\label{lem:monotone}
$\mathrm{DEC}(q)\Rightarrow\mathrm{DEC}(q+1)$ for $q<n$.
\end{lem}

\begin{proof}
Adjoining a zero row to $R$ adds a zero row to $L$, to $LA$ and to $\lambda L-LA$, so it leaves $H$,
$H_A$ and the ranks in \eqref{eq:5a} and \eqref{eq:5b} unchanged; the augmented $R$ is admissible because
$\mathrm{DEC}(q+1)$ does not restrict the free rows.
\end{proof}

\begin{remark}[Coordinate change to $S$]\label{rem:YS}
Fix a rational basis of $S$ and extend it to a rational basis of $\mathbb R^k$; let $Y_S$ be
the change-of-basis matrix, so that $X\mapsto XY_S$ puts the coordinates of $S$ first. Right
multiplication by $Y_S$ preserves ranks and the inclusions and equalities between row spaces,
though not row spaces themselves, so Lemma~\ref{lem:W-conditions} and
Lemma~\ref{lem:struct} apply verbatim in the new coordinates and every object involved lives
in $\mathbb R^\ell$. The basis is obtained by column elimination on $AN$, so the entries of
$Y_S$ have bit length $\operatorname{poly}(\mathcal L)$.
\end{remark}

\section{The finite certificate}\label{sec:certificate}

This section derives the finite certificate for the spectral condition \eqref{eq:5b} in three steps:
a pivot factorization of the rank-$\omega$ block, a rank-consistency test for the linear system that
the factorization leaves, and a compression of that test to a Gram matrix whose size is governed by the
parameter $\ell$ of \eqref{eq:ell}.

Two notions of rank have to be kept apart throughout. The quantity $\omega=\operatorname{rank}H$ is the
\emph{generic} rank of $P(\lambda)=\lambda H-H_A$ over the rational function field $\mathbb R(\lambda)$,
and $\operatorname{rank}P(\lambda)\le\omega$ for every $\lambda$ because $P$ factors through the
rank-$\omega$ matrix $B$. Condition \eqref{eq:5b} asks for the \emph{pointwise} rank to equal $\omega$ at
every $\lambda\in\mathbb C$, that is, for the pointwise rank not to drop below the generic one anywhere.
The certificate below therefore has to express "no drop at any finite $\lambda$", not merely "generic
rank $\omega$"; with $P=(\lambda X-G)B$ this is the constant-rank condition on the pencil $\lambda X-G$.

\subsection{Rank factorization and consistency}\label{sec:consistency}

\paragraph*{Pivot factorization}\label{sec:pivot}
Let $\omega>0$. For a rank-$\omega$ matrix $H$ choose a pivot index set
$J_{p}\subseteq[k]$, $|J_{p}|=\omega$, such that
\begin{equation*}
X=H_{:,J_{p}}\in\mathbb R^{q\times\omega}
\end{equation*}
has full column rank. Then
\begin{equation*}
H=XB,\qquad B\in\mathbb R^{\omega\times k},\qquad B_{:,J_p}=I_\omega .
\end{equation*}
If $H_A=FH$, put $G=FX\in\mathbb R^{q\times\omega}$.
Then
\begin{equation*}
H_A=GB,\qquad P(\lambda)=(\lambda X-G)B .
\end{equation*}
Conversely, if $H_A=GB$ and $X_LX=I_\omega$, where $X_L$ denotes a left
inverse of $X$, then $F=GX_L$ satisfies
\begin{equation*}
FH=GX_LXB=GB=H_A .
\end{equation*}

\paragraph*{Parametric linear system}\label{sec:param}
Let $D$ be a rational matrix with $\ker D=\mathcal T$. For any fixed row index 
$i$, the parameter matrix $B$ is subject to the target constraints
\begin{equation*}
X_iB=h_i^0\in\mathbb R^{1\times k},\qquad G_iB=g_i^0\in\mathbb R^{1\times k},
\end{equation*}
as well as the structural constraint
\begin{equation*}
D\begin{bmatrix}(X_iB)^{\mathsf T}\\ (G_iB)^{\mathsf T}\end{bmatrix}=0.
\end{equation*}
which ensures that the generated row pairs belong to the attainable subspace $\mathcal T$.

Given the pivot factorization condition $B_{:,J_p}=I_\omega$, the known entries of 
$B$ can be eliminated. Let $b$ denote the vectorization of the remaining unknown entries of $B$. The above constraints then reduce to a standard linear system:


\begin{equation*}
\mathcal M(X,G)\,b=c(X,G).
\end{equation*}

Table~\ref{tab:certdim} records the sizes of the objects that enter the certificate; the last row is
the reason the certificate is governed by $\omega$ rather than by the number of maximal minors.

\begin{table}[!t]
\centering
\caption{Sizes of the objects in the certificate of \eqref{eq:5b}. The $\omega\times\omega$ minors are
never formed: their coefficients enter only through the Gram matrix of the last row.}
\label{tab:certdim}
\footnotesize
\begin{tabular}{@{}lll@{}}
\toprule
Object & Size & Bound \\
\midrule
pivot block $X$ & $q\times\omega$ & $\omega\le k$ \\
second block $G$ & $q\times\omega$ & $\omega\le k$ \\
certificate vector $x$ & $2\omega$ & $\omega\le k$ \\
window coefficients of the minors & $|\text{monomials}|$ & degree $\le\omega$ in $\lambda$ \\
Gram matrix $Z^{\mathsf T}Z$ & $2\omega\times2\omega$ & independent of $\binom{k}{\omega}$ \\
\bottomrule
\end{tabular}
\end{table}

\begin{lem}[Rank consistency]\label{lem:rankcons}
For real matrices $M,c$,
\begin{equation*}
\exists b:\ Mb=c
\end{equation*}
if and only if
\begin{equation*}
\operatorname{rank}(MM^{\mathsf T})
=\operatorname{rank}(MM^{\mathsf T}+cc^{\mathsf T}).
\end{equation*}
Equivalently, with $\Gamma=MM^{\mathsf T}$ and
$\Gamma'=MM^{\mathsf T}+cc^{\mathsf T}$,
\begin{equation*}
\bigvee_{t=0}^{d_1}
\left[e_t(\Gamma)>0\wedge e_{t+1}(\Gamma)=0\wedge e_{t+1}(\Gamma')=0\right],
\end{equation*}
where $d_1=\min(\text{rows},\text{cols})$ of $M$, and $e_j$ is the
$j$-th elementary symmetric polynomial of the eigenvalues, with the convention $e_j=0$ whenever $j$
exceeds the size of the matrix.
\end{lem}

\begin{proof}
By the rank criterion for linear systems \cite[Ch.~0]{hj13}, $Mb=c$ has a solution if and only if
$\operatorname{rank}[M\ c]=\operatorname{rank}M$. Since
\begin{equation*}
\begin{gathered}
\operatorname{im}[M\ c]=\operatorname{im}M+\operatorname{span}\{c\},\qquad
\operatorname{im}(MM^{\mathsf T})=\operatorname{im}M,\\
\operatorname{im}(MM^{\mathsf T}+cc^{\mathsf T})=\operatorname{im}[M\ c],
\end{gathered}
\end{equation*}
the two rank equalities of the statement describe the same condition. For a positive semidefinite
$\Gamma$ the elementary symmetric functions of its eigenvalues satisfy
$e_j(\Gamma)>0\iff\operatorname{rank}\Gamma\ge j$ \cite[Ch.~4]{hj13}, so the displayed disjunction
holds exactly when $t=\operatorname{rank}(MM^{\mathsf T})$, and
$e_{t+1}(\Gamma')=0$ together with $\Gamma'\succeq\Gamma$ forces
$\operatorname{rank}(MM^{\mathsf T}+cc^{\mathsf T})=t$.
\end{proof}

\paragraph*{Spectral reduction}\label{sec:spectral}
Since $B$ has full row rank,
\begin{equation*}
\operatorname{rank}P(\lambda)
=\operatorname{rank}\bigl((\lambda X-G)B\bigr)
=\operatorname{rank}(\lambda X-G).
\end{equation*}
Thus \eqref{eq:5b} is equivalent to
\begin{equation*}
\operatorname{rank}(\lambda X-G)=\omega,\qquad\forall\lambda\in\mathbb C .
\end{equation*}

\subsection{The spectral certificate}\label{sec:spectrum}

Two features of the following lemma matter for what comes after it. The hypothesis is on
$f_1,\dots,f_s\in\mathbb R[\lambda]$, but the conclusion gives multipliers $g_i$ that are again
\emph{real} and of degree $<\omega$. This is not a cosmetic refinement: the identity
$\sum_ig_if_i=1$ is then substituted into a system of polynomial equations with real coefficients and
solved over $\mathbb R$, so complex multipliers of bounded degree would not yield a real algebraic
certificate at all. The proof achieves it by choosing the auxiliary combination $b$ among the real
vectors that avoid finitely many hyperplanes, and by invoking the two-polynomial identity over
$\mathbb R[\lambda]$.

\begin{lem}[B\'ezout degree bound]\label{lem:bezout}\label{sec:bezout}
Let $\omega\ge1$, let $s\ge1$, and let $f_1,\ldots,f_s\in\mathbb R[\lambda]$ have
degree at most $\omega$ and no common complex root. Then there exist
real polynomials $g_i$ with $\deg g_i<\omega$ such that
\begin{equation*}
\sum_i g_i f_i=1 .
\end{equation*}
\end{lem}

\begin{proof}
Choose $a\in\mathbb R^s$ with $g=\sum_i a_if_i\neq0$, possible because the $f_i$ have no common
complex root, and let $\alpha_1,\dots,\alpha_m$ be the distinct roots of $g$. For $b\in\mathbb R^s$ put
$h=\sum_i b_if_i$ and write $h(\alpha_j)=\ell_j(b)+i\,m_j(b)$, with $\ell_j,m_j$ real linear forms.
Some $f_i(\alpha_j)$ is nonzero, since otherwise $\alpha_j$ would be a common root, so $(\ell_j,m_j)$
does not vanish identically and $V_j=\{b:h(\alpha_j)=0\}=\ker\ell_j\cap\ker m_j$ is a proper subspace
of $\mathbb R^s$. A finite union of proper subspaces cannot cover $\mathbb R^s$, so some $b$ satisfies
$h(\alpha_j)\neq0$ for every $j$; then $g$ and $h$ are coprime, and the two-polynomial B\'ezout
identity with degree bounds \cite[Ch.~3]{vzgg13} gives $u,v$ with $ug+vh=1$, $\deg u<\deg h\le\omega$
and $\deg v<\deg g\le\omega$. Expanding $g,h$ in the $f_i$ and setting $g_i=ua_i+vb_i$ gives
$\sum_i g_if_i=1$ with $\deg g_i<\omega$.
\end{proof}

\begin{lem}[Degree-windowed relations]\label{lem:window}
Let $\omega\ge1$. Let $f_1,\dots,f_s\in\mathbb R[\lambda]$ have degree at most $\omega$,
let $c_{i,a}$ be the coefficient of $\lambda^a$ (with $c_{i,a}=0$ for
$a<0$ or $a>\omega$), and let $Z$ be the matrix with rows $(i,j)$,
$0\le j<\omega$, columns $t=0,\dots,2\omega-1$, and
$Z_{(i,j),t}=c_{i,t-j}$. For a polynomial $p$ of degree at most
$2\omega-1$, let $\operatorname{coef}(p)\in\mathbb R^{2\omega}$ be its
coefficient vector in the basis $1,\lambda,\dots,\lambda^{2\omega-1}$.
Then
\begin{equation*}
\operatorname{row}Z=
\Bigl\{
\operatorname{coef}\Bigl(\textstyle\sum_i g_if_i\Bigr)
:\ \deg g_i<\omega
\Bigr\}.
\end{equation*}
\end{lem}

\begin{proof}
Row $(i,j)$ of $Z$ is $\operatorname{coef}(\lambda^jf_i)$, because
\begin{equation*}
\lambda^jf_i=\sum_t c_{i,t-j}\lambda^t .
\end{equation*}
For weights $y_{i,j}$, the linearity of $\operatorname{coef}$ on
polynomials of degree at most $2\omega-1$ gives
\begin{equation*}
\begin{aligned}
\sum_{i,j}y_{i,j}\operatorname{coef}(\lambda^jf_i)
&=\operatorname{coef}\Bigl(\sum_{i,j}y_{i,j}\lambda^jf_i\Bigr)\\
&=\operatorname{coef}\Bigl(\sum_i g_if_i\Bigr),
\qquad g_i=\sum_j y_{i,j}\lambda^j,
\end{aligned}
\end{equation*}
with $\deg g_i<\omega$; conversely, every family $g_i$ with
$\deg g_i<\omega$ arises from its own coefficients $y_{i,j}$. Since
$\deg(\lambda^jf_i)\le\omega+j\le2\omega-1$, all coefficients fit in
the $2\omega$ columns, and the convention $c_{i,a}=0$ outside
$0\le a\le\omega$ makes the truncation automatic.
\end{proof}

\begin{prop}[Gram certificate]\label{prop:gram}\label{sec:gram}
Let $\omega=\operatorname{rank}H\ge1$, and assume \eqref{eq:5a}, so that
$P(\lambda)=(\lambda I-F)H$ by Lemma~\ref{lem:compression}(ii). Let $f_i$ be the maximal minors of
$P(\lambda)$, and let $Z$ be the window
matrix of Lemma~\ref{lem:window} built from them, with $c_{i,a}$ the coefficient of $\lambda^a$
in $f_i$; let $e_0=(1,0,\ldots,0)^{\mathsf T}\in\mathbb R^{2\omega}$. Then
\begin{equation*}
\eqref{eq:5b}
\iff e_0^{\mathsf T}\in\operatorname{row}Z
\iff\exists x\in\mathbb R^{2\omega}:\ Z^{\mathsf T}Zx=e_0 .
\end{equation*}
\end{prop}

\begin{proof}
By construction, $\operatorname{coef}(\lambda^jf_i)$ is row $(i,j)$ of $Z$, so
Lemma~\ref{lem:window} identifies $\operatorname{row}Z$ with the coefficient vectors of the sums
$\sum_i g_if_i$ with $\deg g_i<\omega$. With $e_0^{\mathsf T}=\operatorname{coef}(1)$, this gives
$e_0^{\mathsf T}\in\operatorname{row}Z$ if and only if there are $g_i$ with $\deg g_i<\omega$ and
$\sum_i g_if_i=1$, which is the B\'ezout representation of Lemma~\ref{lem:bezout}; it exists if
and only if the $f_i$ have no common complex root, that is, if and only if \eqref{eq:5b}. The
window $\deg g_i<\omega$ is exactly the one Lemma~\ref{lem:bezout} supplies, so no bound on an
arbitrary relation is needed. Finally, for real $Z$,
\begin{equation*}
\operatorname{row}Z=\operatorname{im}Z^{\mathsf T}
=\operatorname{im}(Z^{\mathsf T}Z),
\end{equation*}
the last equality because $\ker Z=\ker(Z^{\mathsf T}Z)$; hence
$e_0^{\mathsf T}\in\operatorname{row}Z$ if and only if
$Z^{\mathsf T}Zx=e_0$ has a real solution.
\end{proof}

\paragraph*{Cauchy--Binet compression}\label{sec:cb}
Work in the coordinates of Remark~\ref{rem:YS}, in which the first $\ell$ vectors are a basis of
$S=\operatorname{row}(AN)$ and the remaining $k-\ell$ ones span a complement. For a feasible $W$ the
two projections coincide and lie in $S$ (Section~\ref{sec:US}), so every row of $H=LN$ and of
$H_A=LAN$ belongs to $S$; hence in these coordinates the last $k-\ell$ columns of $P(\lambda)=\lambda
H-H_A$ vanish, and a nonzero minor uses only the $\ell$ columns of $S$. Accordingly define
\begin{equation*}
K(\lambda,\mu)=\sum_{|J|=\omega}
\det\left(P(\lambda)_{:,J}^{\mathsf T}P(\mu)_{:,J}\right),
\end{equation*}
where $J$ ranges over the $\omega$-element subsets of $[\ell]$; a column set meeting the
complement of $S$ gives a zero determinant. For each fixed column set $J$, the Cauchy--Binet identity
\cite[Ch.~0]{hj13} applied
to the $\omega\times q$ matrix $P(\lambda)_{:,J}^{\mathsf T}$ and the
$q\times\omega$ matrix $P(\mu)_{:,J}$ gives
\begin{equation*}
\det\!\left(P(\lambda)_{:,J}^{\mathsf T}P(\mu)_{:,J}\right)
=\sum_{\substack{I\subseteq[q]\\ |I|=\omega}}
f_{I,J}(\lambda)f_{I,J}(\mu),
\end{equation*}
where $f_{I,J}$ is the minor on rows $I$ and columns $J$. Summing over $J$,
\begin{equation*}
K(\lambda,\mu)=\sum_{\substack{I\subseteq[q]\\ |I|=\omega}}
\sum_{\substack{J\subseteq[\ell]\\ |J|=\omega}}
f_{I,J}(\lambda)f_{I,J}(\mu),
\end{equation*}
the Gram kernel of \textbf{all} maximal minors of $P$.

Three consequences are used below.
\begin{enumerate}[label=(\roman*)]
\item\label{item:cost} \emph{Cost.} Computing $K$ needs only
$\binom{\ell}{\omega}\le2^{\ell}$ determinants of order $\omega$, one per
column set $J$ of the $\ell$-dimensional space $S$ of \eqref{eq:ell}; the
$\binom q\omega\binom k\omega$ individual minors are never enumerated,
and no hidden enumeration replaces them.
\item\label{item:coef} \emph{Coefficient identity.} Writing $S_{a,b}$
for the coefficient of $\lambda^a\mu^b$ in $K$ and $c_{i,a}$ for the
coefficient of $\lambda^a$ in $f_i$, one has
$S_{a,b}=\sum_i c_{i,a}c_{i,b}$. Hence, with the row indexing of
Lemma~\ref{lem:window},
\begin{equation*}
(Z^{\mathsf T}Z)_{t,t'}
=\sum_{j=0}^{\omega-1}\sum_i c_{i,t-j}c_{i,t'-j}
=\sum_{j=0}^{\omega-1}S_{t-j,t'-j}.
\end{equation*}
\item\label{item:roots} \emph{Common roots.} Since the coefficients of
$P$ are real, specializing at $(\lambda,\mu)=(\alpha,\bar\alpha)$ gives
\begin{equation*}
K(\alpha,\bar\alpha)
=\sum_{I,J}\left|f_{I,J}(\alpha)\right|^2,
\end{equation*}
so $K(\alpha,\bar\alpha)=0$ if and only if every maximal minor of $P$
vanishes at $\alpha$. This identity shows that the coefficient matrix of
$K$ encodes precisely the common-root condition; the algorithm below
nevertheless decides that condition at the coefficient level, without
evaluating $K$ at complex points.
\end{enumerate}

\begin{remark}[Algorithmic reading]\label{rem:algo}
The compression is computed in four steps:
(i) for each column set $J\subseteq[k]$, $|J|=\omega$, form the
$\omega\times\omega$ pencil $P(\lambda)_{:,J}^{\mathsf T}P(\mu)_{:,J}$;
(ii) compute its bivariate coefficients, of degree at most $\omega$ in
each variable;
(iii) sum the coefficients over $J$;
(iv) recover $Z^{\mathsf T}Z$ from the shifted identity of
item~\ref{item:coef}.
The identity is classical (an exterior-power Gram identity; see, e.g.,
\cite{emi99} for the B\'ezoutian/elimination background). What the
construction adds is the combination with the parameter-sized
certificate: by the cost bound of item~\ref{item:cost}, the certificate
matrix stays parameter-sized when $q$ is large and $k$ is bounded.
\end{remark}

\section{The fixed-order algorithm}\label{sec:fixedq}

This section decides $\mathrm{DEC}(q)$ at a prescribed order $q$. The procedure
is Algorithm~\ref{alg:fixed-q}: in each pivot chart it quantifies only over the
entries of $X$, $G$ and the certificate vector, that is over
$2q\omega+2\omega=O(q^2)$ real variables, whereas the chart coordinates of the
order-by-order procedure of \cite{fer26b} number $d(m_L-d)=(n-p-q)(q-r)$ and
still grow linearly in $n$ when $q$ is fixed; the compression of
Section~\ref{sec:compressed} is what removes that dependence.

The inner test of Algorithm~\ref{alg:fixed-q} is the only step that is not
elementary, and it is worth stating separately. For a branch $(\omega,J)$ the unknown is a triple
$(X,G,x)$ of real matrices, and the branch is the system
\begin{equation}\label{eq:phi}
\Phi_{\omega,J}:\qquad
\begin{cases}
\text{(i)}\ \ e_\omega(X^{\mathsf T}X)>0,\\
\text{(ii)}\ \ \mathcal M(X,G)\,b=c(X,G)\ \text{is consistent},\\
\text{(iii)}\ \ X_iB=h_i^0,\quad G_iB=g_i^0\qquad(i\le r_0),\\
\text{(iv)}\ \ D\,[(X_iB)^{\mathsf T};(G_iB)^{\mathsf T}]=0\qquad(i>r_0),\\
\text{(v)}\ \ Z^{\mathsf T}Zx=e_0,
\end{cases}
\end{equation}
with $B$ and $Z$ as in Section~\ref{sec:certificate}. Items (i)--(iv) are equations and one strict
inequality over the reals; (v) is the Gram form of the spectral condition, also polynomial. Asking
whether $\Phi_{\omega,J}$ has a real solution is therefore a decision problem in the existential theory
of the reals, and it is decided by the real-algebraic procedure of Appendix~\ref{app:bits}, whose cost
in the numbers of variables, polynomials, degree and coefficient bit length is what step (ii) of the
proof of Theorem~\ref{thm:fixedq} counts. Nothing is enumerated except the discrete pair $(\omega,J)$.

\begin{algorithm}[!t]
\caption{Deciding $\mathrm{DEC}(q)$ for a fixed order $q$}
\label{alg:fixed-q}
\begin{algorithmic}[1]
\Require $A,C,L_0,q$; \quad $r_0\gets\operatorname{rows}(L_0)$; \quad $r\gets q-r_0$
\Ensure \textsc{feasible} or \textsc{infeasible}
\State $N\gets$ basis of $\ker[C;CA]$
\State $h^0\gets L_0N$;\quad $g^0\gets L_0AN$
\State $D\gets$ basis of $\bigl(\operatorname{row}[N\ AN]\bigr)^{\perp}$
\If{$r=0$}
  \State \Return the exact test of Section~\ref{sec:zero}
\EndIf
\If{$h^0=0$ \textbf{and} $g^0=0$}
  \State \Return \textsc{feasible}
  \Statex \Comment{$W_0=0$, so $W=0$ is feasible}
\EndIf
\For{$\omega=1$ \textbf{to} $\min(q,k)$}
  \For{each $J\subseteq[k]$ with $|J|=\omega$}
    \If{$\Phi_{\omega,J}$ of \eqref{eq:phi} has a real solution}
      \State \Return \textsc{feasible}
    \EndIf
  \EndFor
\EndFor
\State \Return \textsc{infeasible}
\end{algorithmic}
\end{algorithm}

\begin{thm}\label{thm:fixedq}
For every fixed $q\ge1$, $\mathrm{DEC}$ is decidable in deterministic polynomial time: with
$r=q-r_0$, Algorithm~\ref{alg:fixed-q} runs in $\mathcal L^{O(q^2)}$ bit operations.
\end{thm}

\begin{proof}
It suffices to show that $\Phi_{\omega,J}$ is satisfiable exactly when a feasible $L$ exists, and then
to substitute the sizes of $\Phi$ into \eqref{eq:rengar} chart by chart.
\emph{(i) Equivalence.} Let $L=[L_0;R]$ be feasible and let $X,B,G$ be as in the pivot
factorization of Section~\ref{sec:pivot}, with $H=LN$, $H_A=LAN$, $\omega=\operatorname{rank}H$ and
$H_A=GB$ from Lemma~\ref{lem:compression}(ii). Then $X_iB=h_i^0$, $G_iB=g_i^0$ for $i\le r_0$ and
$D[(X_iB)^{\mathsf T};(G_iB)^{\mathsf T}]=0$ for $i>r_0$, by the definition of $H$ and of $D$; the
remaining three conditions of $\Phi_{\omega,J}$ are $e_\omega(X^{\mathsf T}X)>0$
($\operatorname{rank}X=\omega$), the rank-consistency condition of Lemma~\ref{lem:rankcons}, and
$Z^{\mathsf T}Zx=e_0$ for \eqref{eq:5b} by Proposition~\ref{prop:gram}. Conversely, if $\Phi_{\omega,J}$
holds, Lemma~\ref{lem:rankcons} gives $B$ with $B_{:,J}=I_\omega$, and $H=XB$, $H_A=GB$ reproduce the
prescribed rows and place the others in $\mathcal T$, so that $H=LN$ and $H_A=LAN$ for $L=[L_0;R]$ with
those rows; since $X$ has a left inverse $X_L$, $H_A=GB=(GX_L)XB$, which is \eqref{eq:5a} by
Lemma~\ref{lem:compression}(i), and $Z^{\mathsf T}Zx=e_0$ is \eqref{eq:5b} by
Proposition~\ref{prop:gram}. Hence $\bigvee_{\omega,J}\Phi_{\omega,J}$ holds exactly when a feasible $L$
exists, which is the test performed by Algorithm~\ref{alg:fixed-q}.

\emph{(ii) Count.} With $\omega\le k$ and $N_{\rm var}=2q\omega+2\omega$ real variables,
\begin{align*}
N &= 2q\omega+2\omega\le 2q^2+2q,\\
s &= O(k)\cdot O(k)+O(k)=O(k^2),\\
D &= \max\{O(k),\,2\omega+1\}=O(k)\le O(n),\\
\beta &= \operatorname{poly}(n,\mathcal L),
\end{align*}
where the two factors of $s$ count the rank-consistency branches
$t\le\min\{\text{rows},\text{cols}\}(M)$, at most $O(k)$ of them, and the
$O(k)$ polynomials of degree $O(k)$ carried by each branch, while $2\omega+1$
is the degree of the Gram and certificate constraints; the last line is the
coefficient-size lemma of Appendix~\ref{app:bits}. Renegar's Part II bound
$\beta(\log\beta)(\log\log\beta)(sD)^{O(N_{\rm var})}$ therefore gives, per chart,
\begin{equation*}
\beta(\log \beta)(\log\log \beta)(sD)^{O(N_{\rm var})}
=\operatorname{poly}(n)^{O(q^2)}
=\mathcal L^{O(q^2)},
\end{equation*}
because $sD=\operatorname{poly}(n)$ and $N\le 2q^2+2q$ with $q$ fixed. The
charts number
\begin{equation*}
N_{\mathrm{ch}}=(q+1)k^{q}\le 2^{q}k^{q}\le\mathcal L^{2q},
\end{equation*}
since $k\le n\le\mathcal L$ and $q+1\le 2^{q}\le\mathcal L^{q}$ for $q\ge1$
and $\mathcal L\ge2$. Hence
\begin{equation*}
N_{\mathrm{ch}}\cdot\mathcal L^{O(q^2)}
\le\mathcal L^{2q}\cdot\mathcal L^{O(q^2)}
=\mathcal L^{O(q^2)} .
\end{equation*}
For fixed $q$ the exponent $O(q^2)$ is a constant, so the construction size
$k^{O(q^2)}\operatorname{poly}(\mathcal L)=\mathcal L^{O(q^2)}$ matches the bound above. The
degenerate cases $k=0$,
$\omega=0$ and $r=0$ are treated in Section~\ref{sec:deg}.
\end{proof}



Three representation details are worth recording.

\begin{enumerate}
\item \emph{Strict inequalities.} A condition $f>0$ is written as
\begin{equation*}
f>0\iff\exists u:\ fu^2=1,
\end{equation*}
which is equivalent over the reals. There are $O(1)$ strict inequalities per
branch, so this adds $O(1)$ variables and changes no order.
\item \emph{Disjunctions.} The rank-consistency disjunction is processed branch
by branch, with one real-feasibility query per $t$ and the running times
added, $t\cdot\mathcal L^{O(q^2)}=O(k)\,\mathcal L^{O(q^2)}=\mathcal L^{O(q^2)}$.
This is why $t$ appears in $s$ and not in $D$.
\item \emph{Elementary symmetric polynomials.} $e_j(G)$ is represented by the
coefficients of $\det(zI+G)$, computed by fraction-free elimination on the
$k\times k$ polynomial matrix $zI+G$ in $O(k^3)$ polynomial operations of
degree $O(k)$. The expansion is never written out term by term.
\end{enumerate}

\paragraph*{Degenerate cases}\label{sec:deg}

\begin{enumerate}
\item \emph{$k=0$}, that is, $\ker[C;CA]=0$. Then $N$ has no columns, so $H=LN=0$ and
$H_A=LAN=0$ for every $L$; both \eqref{eq:5a} and \eqref{eq:5b} hold, for $L=L_0$ in particular, so
$q_{\min}=r_0$.
\item \emph{$\omega=0$ with $k>0$.} Then $H=0$, and
$\operatorname{rank}(\lambda H-H_A)=0$ for all $\lambda\in\mathbb C$ if and only if $H_A=0$,
which together with $H=0$ is exactly $\eqref{eq:5a}\wedge\eqref{eq:5b}$ for
that $L$. The branch $H=H_A=0$ of the encoding is therefore neither stronger
nor weaker than the condition.
\item \emph{$q<r_0$.} Rejected by definition, since $L=[L_0;R]$ has $q$ rows
and contains the $r_0$ rows of $L_0$.
\item \emph{$r=0$.} Decided by the exact pencil test of Section~\ref{sec:zero}.
\end{enumerate}

\section{The fixed-parameter algorithm}\label{sec:fpt}

By Lemma~\ref{lem:sandwich} a feasible observer corresponds to a subspace $W$ with
$W_0\subseteq W\subseteq\mathcal T$ and $\dim(W/W_0)\le r$, and by Lemma~\ref{lem:W-conditions} the
two conditions depend on $W$ alone. Lemma~\ref{lem:struct} confines the feasible $W$ to
$\omega+1\le\dim W\le2\omega$, Section~\ref{sec:US} bounds $\omega$ by $\ell=\dim S$, and
Remark~\ref{rem:YS} puts every object in
$\mathbb R^\ell$. Algorithm~\ref{alg:fpt} does this: it quantifies over a chart of a Grassmannian of dimension at most
$2\ell$ and over the $2\omega$ certificate variables, and the only non-polynomial factor
is the parameter $\ell=\operatorname{rank}(AN)$, which by \eqref{eq:ell} does not involve $r$.

\paragraph*{Charts and their coordinates}\label{sec:chartcoords}
The charts of Algorithm~\ref{alg:fpt} are written with two fixed matrices. With $W_0$ and
$\mathcal V$ as in the algorithm, put $s_0=\dim W_0$ and $m'=\dim\mathcal V-s_0$. Let
$E_0\in\mathbb R^{s_0\times2k}$ be a matrix whose rows are a basis of $W_0$, and let
$E_Q\in\mathbb R^{m'\times2k}$ be a matrix whose rows complete that basis to a basis of $\mathcal V$;
both have rational entries and are computed once per instance. The point of the chart of
$\mathrm{Gr}(t,m')$ indexed by $Y\in\mathbb R^{t\times(m'-t)}$ is the row space of
\begin{equation*}
Z_{\rm ch}(Y)=[E_0;\,[I_t\mid Y]E_Q]\in\mathbb R^{(s_0+t)\times2k},
\end{equation*}
which contains $W_0$ by construction. In the coordinates of Remark~\ref{rem:YS}, in which every row
of $H$ and $H_A$ lies in $S$ (Section~\ref{sec:US}), the two blocks of $\ell$ columns of this matrix
are the matrices called $H(Y)$ and $G(Y)$ below: the \emph{complete} $H=LN$ and $H_A=LAN$, that is,
including the $s_0$ rows contributed by $L_0$, so that \eqref{eq:5a} is tested on all rows of $H$ and
$H_A$ and not on the $t$ free rows alone.

\begin{algorithm}[!t]
\caption{Deciding $\mathrm{DEC}(q)$ in time $f(\ell)\operatorname{poly}(\mathcal L)$}
\label{alg:fpt}
\begin{algorithmic}[1]
\Require $A,C,L_0,q$ as above; $r=q-r_0$.
\Ensure \textsc{feasible} or \textsc{infeasible}.
\If{$r=0$}
  \State \Return the exact test of Proposition~\ref{prop:zero} on $L=L_0$
\EndIf
\If{$q<q_0$}
  \State \Return \textsc{infeasible} \Comment{the algebraic lower bound \eqref{eq:range}}
\EndIf
\State $N\gets$ basis of $\ker[C;CA]$;\quad $M\gets[N\ AN]$;\quad
$S\gets\operatorname{row}(AN)$;\quad $\ell\gets\operatorname{rank}(AN)$
\State $W_0\gets\operatorname{row}(L_0M)$;\quad
$\mathcal V\gets\mathcal T\cap(S\times S)$
\If{$W_0\not\subseteq\mathcal V$}
  \State \Return \textsc{infeasible} \Comment{necessary: every feasible $W$ satisfies $W_0\subseteq W\subseteq S\times S$}
\EndIf
\If{$W_0=0$}
  \State \Return \textsc{feasible} \Comment{the $\omega=0$ branch: $W=0$ is realized by $R=0$}
\EndIf
\State $s_0\gets\dim W_0$;\quad $m'\gets\dim\mathcal V-s_0$
\For{$t=0$ \textbf{to} $\min(r,2\ell)$}
  \For{each chart of $\mathrm{Gr}(t,m')$ with chart coordinates $Y\in\mathbb R^{t\times(m'-t)}$}
    \State $H(Y),G(Y)\gets$ the two $(s_0+t)\times\ell$ blocks of the chart matrix
             $Z_{\rm ch}(Y)=[E_0;\,[I_t\mid Y]E_Q]$; write $H,G$ for them
    \Statex \quad \Comment{$H=LN$ and $G=LAN$, including the $s_0$ rows of $L_0$}
    \For{$\omega=1$ \textbf{to} $\ell$}
      \State form $\Psi_{t,\omega}(Y,x)$ with
      \Statex \quad (i) $e_\omega(HH^{\mathsf T})>0$ and
               $e_{\omega+1}\bigl([H;G][H;G]^{\mathsf T}\bigr)=0$;
      \Statex \quad (ii) $Z(Y)^{\mathsf T}Z(Y)x=e_0$, with $Z(Y)$ the window matrix of
               Lemma~\ref{lem:window} for $\lambda H-G$, computed from $K(\lambda,\mu)$ by
               Cauchy--Binet over the $\ell$ columns of $S$.
      \If{$\Psi_{t,\omega}$ is real-feasible}
        \State \Return \textsc{feasible}
      \EndIf
    \EndFor
  \EndFor
\EndFor
\State \Return \textsc{infeasible}
\end{algorithmic}
\end{algorithm}

\begin{thm}\label{thm:fpt}
For rational input and every $r\ge0$, $\mathrm{DEC}$ is decidable in
\begin{equation}\label{eq:fpt-l}
\exp\bigl(O(\ell^2\log(\ell+2))\bigr)\operatorname{poly}(\mathcal L)
\end{equation}
bit operations, where $\ell=\operatorname{rank}(AN)\le k$ does not depend on the number $r$ of
free rows. For $r=0$ the bound sharpens to $\operatorname{poly}(\mathcal L)$.
\end{thm}

\begin{proof}
The case $r=0$ is Proposition~\ref{prop:zero}. For $r\ge1$ it suffices to parameterize every feasible
$W$ by the chart coordinates of $W/W_0\subseteq\mathcal V/W_0$ and to express the two conditions on
that chart by a formula whose size is governed by $\ell$; the argument has six steps.

\emph{(i) The bounds $\omega\le\ell$ and $t\le2\ell$.} By Section~\ref{sec:US} a feasible $W$ has
$\pi_1(W)=\pi_2(W)=U\subseteq S$, so $\omega=\dim U\le\dim S=\ell$; by Lemma~\ref{lem:struct}
$\dim W\le2\omega\le2\ell$, so the quotient dimension $t=\dim(W/W_0)\le\dim W\le2\ell$ as well.
These are the ranges $t\le\min(r,2\ell)$ and $\omega\le\ell$ of the algorithm.

\emph{(ii) Reduction to a subspace.} By Lemma~\ref{lem:sandwich} a feasible free block exists if and
only if some feasible $W$ satisfies $W_0\subseteq W\subseteq\mathcal T$ and
$\dim(W/W_0)\le r$, and by Lemma~\ref{lem:monotone} the bound $r$ can be saturated; it is therefore
enough to range $t=\dim(W/W_0)$ over $0,\dots,\min(r,2\ell)$ and to test each $t$ once.

\emph{(iii) The inclusion test is necessary.} If $\pi_1(W)=\pi_2(W)=U\subseteq S$
(Section~\ref{sec:US}), then $W\subseteq U\times U\subseteq S\times S$, and $W\subseteq\mathcal T$
by Lemma~\ref{lem:sandwich}; hence $W_0\subseteq W\subseteq\mathcal V=\mathcal T\cap(S\times S)$.
The test at the head of the algorithm therefore rejects no feasible instance, and it is what makes
$m'=\dim\mathcal V-s_0$ a nonnegative number.

\emph{(iv) The charts exhaust the candidates.} Every $W$ with
$W_0\subseteq W\subseteq\mathcal V$ and $\dim(W/W_0)=t$ maps to the subspace $W/W_0$ of
$\mathcal V/W_0$ of dimension $t$, hence lies in one of the charts of $\mathrm{Gr}(t,m')$ that the
algorithm enumerates.

\emph{(v) The formula characterizes the two conditions.} Within a chart, $\Psi_{t,\omega}$ is
satisfiable exactly when the blocks $(H(Y),G(Y))$ of the chart matrix $Z_{\rm ch}(Y)$ satisfy
\eqref{eq:5a} with $\operatorname{rank}H(Y)=\omega$ and \eqref{eq:5b}, by
Lemma~\ref{lem:compression} and Lemma~\ref{lem:window}.

Fix $t$. By Lemma~\ref{lem:struct} and Remark~\ref{rem:YS} a feasible $W$ of quotient dimension
$t$ is a point of the Grassmannian $\mathrm{Gr}(t,m')$ with $m'\le2\ell-s_0$, written in the chart
coordinates $Y\in\mathbb R^{t\times(m'-t)}$ as the row space of
$Z_{\rm ch}(Y)=[E_0;[I_t\mid Y]E_Q]$ of Section~\ref{sec:chartcoords}, which
contains $W_0$ by construction. The chart count is
\begin{equation}\label{eq:charts}
\sum_{t}\binom{m'}{t}=2^{m'}\le4^{\ell}.
\end{equation}

Within a chart and for a fixed $\omega$, the quantified variables are $Y$ and the certificate
$x\in\mathbb R^{2\omega}$, and the data of the real-feasibility problem are
\begin{equation}\label{eq:nvar}
\begin{gathered}
N_{\rm var}=t(m'-t)+2\omega\le\tfrac{m'^2}{4}+2\omega\le\ell^2+2\ell,\\
s=2\omega+O(1)\le2\ell+O(1),\qquad D=O(\ell),\\
\beta=\operatorname{poly}(\mathcal L),
\end{gathered}
\end{equation}
where the discrete choices of $t$, of the chart and of $\omega$ are finite disjunctions rather
than real variables, and $\beta$ is the coefficient bit length of Lemma~\ref{lem:coeff}. The
polynomials are those of Lemma~\ref{lem:window}: $e_\omega(H(Y)H(Y)^{\mathsf T})>0$ and
$e_{\omega+1}\bigl([H(Y);G(Y)][H(Y);G(Y)]^{\mathsf T}\bigr)=0$ encode \eqref{eq:5a} together with
$\operatorname{rank}H(Y)=\omega$ (Lemma~\ref{lem:compression}), and
$Z(Y)^{\mathsf T}Z(Y)x=e_0$ encodes \eqref{eq:5b}. Substituting \eqref{eq:nvar} into \eqref{eq:rengar} gives, per
chart,
\begin{equation*}
\begin{gathered}
T\le\beta(\log\beta)(\log\log\beta)(sD)^{O(N_{\rm var})}\\
=\operatorname{poly}(\mathcal L)\exp\bigl(O(\ell^2\log(\ell+2))\bigr),
\end{gathered}
\end{equation*}
and multiplying by the chart count \eqref{eq:charts} absorbs $\exp(O(\ell))$ into the same
exponential. The total is \eqref{eq:fpt-l}; Table~\ref{tab:fptcount} collects the sources.

\emph{(vi) Conclusion.} The algorithm has four exits, and each is correct. If $r=0$ it returns the
exact test of Section~\ref{sec:zero}, correct by Proposition~\ref{prop:zero}. If $q<q_0$ it returns
\textsc{infeasible}, correct by the algebraic lower bound \eqref{eq:range}, which every feasible $L$
satisfies. If $W_0=0$ it returns \textsc{feasible}, because then $W=0$ is admissible and is realized
by the zero free block. Otherwise $r\ge1$, $q\ge q_0$ and $W_0\ne0$, and steps (i)--(v) show that the
loop over $t$, over the charts and over $\omega$ accepts exactly when a feasible $W$ with
$t=\dim(W/W_0)\le\min(r,2\ell)$ exists. If it does, the converse half of Lemma~\ref{lem:sandwich}
realizes $W$ by a free block with exactly $t$ rows, and Lemma~\ref{lem:monotone}, applied
$q-r_0-t$ times, pads that block with zero rows, so the observer has exactly $q$ rows; if it does
not, Lemma~\ref{lem:sandwich} excludes every feasible $L$. Hence the algorithm accepts if and only
if a feasible observer exists.
\end{proof}

\begin{table}[!t]
\centering
\caption{Sources of the fixed-parameter bound \eqref{eq:fpt-l}, each counted as in the proof.}
\label{tab:fptcount}
\footnotesize
\begin{tabular}{@{}ll@{}}
\toprule
Source & Bound \\
\midrule
quotient dimension $t$ (one per order) & at most $2\ell+1$ \\
Grassmann charts for $t$ & $\sum_t\binom{m'}{t}=2^{m'}\le4^{\ell}$ \\
rank $\omega$ of the compressed block & at most $\ell$ \\
real variables of the chart and certificate & $\ell^2+2\ell$ \\
polynomials per branch & $s=2\omega+O(1)\le2\ell+O(1)$ \\
maximum degree & $O(\ell)$ \\
coefficient bit length & $\operatorname{poly}(\mathcal L)$ \\
\bottomrule
\end{tabular}
\end{table}

\begin{cor}\label{cor:arith}
The same algorithm uses $\exp\bigl(O(\ell^2\log(\ell+2))\bigr)\operatorname{poly}(\mathcal L)$
arithmetic operations.
\end{cor}

\begin{proof}
Immediate from Theorem~\ref{thm:fpt}: arithmetic operations are among the operations counted in the
bit bound, and the data have bit length $\operatorname{poly}(\mathcal L)$.
\end{proof}

\section{The zero-free-row case: an exact polynomial-time pencil test}\label{sec:zero}

For $r=0$ only two conditions remain, and both are decidable without a
parameter-dependent search. \eqref{eq:5a} is a constant-rank test on
$H=L_0N$, $H_A=L_0AN$ (Lemma~\ref{lem:compression}(i)); \eqref{eq:5b} is the
constancy of the rank of the degree-one pencil $\lambda X-G$. We first isolate
the pencil statement.

\begin{lem}\label{lem:zero}
Let $X,G\in\mathbb Q^{q\times\omega}$ with $q\ge\omega\ge1$, and put
$Q(\lambda)=\lambda X-G$. The following are equivalent:
\emph{(i)} $\operatorname{rank}Q(\lambda)=\omega$ for every
$\lambda\in\mathbb C$;
\emph{(ii)} the maximal $\omega\times\omega$ minors of $Q$ have no common
complex root;
\emph{(iii)} the maximal minors do not all vanish and their gcd in
$\mathbb Q[\lambda]$ is a nonzero constant;
\emph{(iv)} there is $S(\lambda)\in\mathbb Q[\lambda]^{\omega\times q}$ with
$S(\lambda)Q(\lambda)=I_\omega$ and $\deg S\le2\omega-2$;
\emph{(v)} there exist matrices
$S_0,\dots,S_{2\omega-2}\in\mathbb Q^{\omega\times q}$ satisfying the rational
linear system
\begin{equation}\label{eq:R0}
\begin{aligned}
-S_0G&=I_\omega,\\
S_{j-1}X-S_jG&=0\quad(1\le j\le2\omega-2),\\
S_{2\omega-2}X&=0 .
\end{aligned}
\end{equation}
\end{lem}

\begin{proof}
The three equivalent forms (i)--(iii) are those of Lemma~\ref{lem:compression}(iii) with $(H,H_A)$
replaced by $(X,G)$ and $k$ by $\omega$, the gcd being taken in $\mathbb Q[\lambda]$ as in
Section~\ref{sec:notation}.

\emph{(iii)$\Rightarrow$(iv).} For each $\omega$-element row set
$I\subseteq[q]$, let $E_I\in\mathbb Q^{\omega\times q}$ be the row selection,
$Q_I=E_IQ$ and $f_I=\det Q_I$. Then
\begin{equation*}
\operatorname{adj}(Q_I)E_IQ=f_II_\omega,\qquad
\sum_Ig_If_I=1\ \text{ with }\ \deg g_I<\omega
\end{equation*}
by Lemma~\ref{lem:bezout}. The coefficients of the $g_I$ solve a linear system with rational data,
because the $f_I$ have rational coefficients; a rational linear system that is solvable over
$\mathbb R$ is solvable over $\mathbb Q$, so the $g_I$ may be taken rational. The cofactors of $Q_I$
are rational polynomials as well, so that
\begin{equation*}
S=\sum_Ig_I\operatorname{adj}(Q_I)E_I
\quad\text{satisfies}\quad
S(\lambda)Q(\lambda)=I_\omega ,
\end{equation*}
and $\deg\operatorname{adj}(Q_I)\le\omega-1$ gives
$\deg S\le(\omega-1)+(\omega-1)=2\omega-2$.

\emph{(iv)$\Rightarrow$(v).} Write $S=\sum_{j=0}^{2\omega-2}S_j\lambda^j$. Comparing the
coefficients of $\lambda^0$, $\lambda^k$ for $1\le k\le2\omega-2$ and $\lambda^{2\omega-1}$ in
$S(\lambda)(\lambda X-G)=I_\omega$ gives exactly \eqref{eq:R0}.

\emph{(v)$\Rightarrow$(i).} Evaluating $S(\lambda)Q(\lambda)=I_\omega$ at $\lambda_0\in\mathbb C$
exhibits a left inverse of $Q(\lambda_0)$, so $\omega\le\operatorname{rank}Q(\lambda_0)\le\omega$,
the upper bound because $Q(\lambda_0)$ has $\omega$ columns.
\end{proof}

The size of \eqref{eq:R0} is fixed by $\deg S\le2\omega-2$: the unknowns are the entries of
$2\omega-1$ matrices of size $\omega\times q$, so $N_{\rm var}=(2\omega-1)q\omega=O(q\omega^2)$,
and the $2\omega$ matrix equations give $N_{\rm eq}=2\omega^3$ scalar equations. All coefficients
are entries of $X$, $G$, or $0,\pm1$, so clearing denominators gives an integer system of size
$M=O(\omega^3+q\omega^2)$ with entries of bit length $O(\mathcal L_0)$, where $\mathcal L_0$ is the
input length of $(X,G)$. Exact fraction-free elimination costs $O(M^3)$ arithmetic operations on
integers of $O(M(\mathcal L_0+\log M))$ bits by Hadamard's inequality, hence
$\operatorname{poly}(q,\omega,\mathcal L_0)$ bit operations, and consistency of \eqref{eq:R0} is
decided within that bound.

\begin{remark}[Construction and algorithm]\label{rem:construction}
The adjugate construction in (iii)$\Rightarrow$(iv) of Lemma~\ref{lem:zero}
enumerates the $\binom q\omega$ maximal minors, which is exponential in $q$; it
is used only to prove the existence of a certificate of degree at most
$2\omega-2$. It applies to the rectangular pencil $Q$ because each $Q_I$ is
square, so that $\deg S\le2\omega-2$ holds for every $q\ge\omega$ with no square
assumption on $Q$. The algorithm never forms the sum: it searches for the
coefficients of $S$ directly in \eqref{eq:R0}, whose size is polynomial in
$q,\omega$. This is the step that turns a certificate-size statement into a
decision algorithm.
\end{remark}

\begin{remark}[Higher degree]\label{rem:degree}
The same argument applies to pencils of degree $d$: if the entries of $Q$ have
degree at most $d$, then every maximal minor has degree at most $d\omega$,
B\'ezout gives multipliers with $\deg g_I<d\omega$, and the adjugate contributes
at most $d(\omega-1)$,
\begin{equation*}
\deg S\le(d\omega-1)+d(\omega-1)=2d\omega-d-1,
\end{equation*}
which is $2\omega-2$ for $d=1$. Degree bounds for polynomial-matrix left
inverses are studied classically \cite{fab90}; the bound above is the elementary
adjugate--B\'ezout one, and only the case $d=1$ is used in this paper. The
general bound may be useful if the observer model is extended. For $d=0$ the
pencil is constant and the statement is the rank test of
Lemma~\ref{lem:compression}(iii).
\end{remark}

\begin{prop}\label{prop:zero}
Let $r=0$, so $L=L_0$. Feasibility of $\eqref{eq:5a}\wedge\eqref{eq:5b}$ is
decidable in $\operatorname{poly}(n,\mathcal L)$ bit operations.
\end{prop}

\begin{proof}
Compute a basis $N$ of $\ker[C;CA]$ and set $H=L_0N$, $H_A=L_0AN$.

\emph{Step 1.} Check $\operatorname{rank}[H_A;H]=\operatorname{rank}H$; if it
fails, reject, since this is \eqref{eq:5a}
(Lemma~\ref{lem:compression}(i)).

\emph{Step 2.} Put $\omega=\operatorname{rank}H$. If $\omega=0$, then $H=0$ and
\eqref{eq:5b} reads $\operatorname{rank}(\lambda H-H_A)=0$ for all $\lambda$,
i.e.\ $H_A=0$; accept exactly in that case.

\emph{Step 3.} Otherwise take a rank factorization $H=XB$ with
$X\in\mathbb Q^{q\times\omega}$ of full column rank and
$B\in\mathbb Q^{\omega\times k}$ of full row rank. Step 1 gives
$\operatorname{row}H_A\subseteq\operatorname{row}H$
(Lemma~\ref{lem:compression}(i)), hence $H_A=FH$ for some $F$; set $G=FX$, so
$H_A=GB$.

\emph{Step 4.} Run the test of Lemma~\ref{lem:zero} on $(X,G)$.

\emph{Correctness.} Since $B$ is of full row rank,
$\operatorname{rank}(\lambda H-H_A)=\operatorname{rank}((\lambda X-G)B)=\operatorname{rank}(\lambda X-G)$
as in Section~\ref{sec:spectral}, for every $\lambda$, and \eqref{eq:5b} is
$\operatorname{rank}(\lambda H-H_A)=\omega$ for all $\lambda$
(Lemma~\ref{lem:compression}(iii)); thus
Lemma~\ref{lem:zero} decides \eqref{eq:5b}. If \eqref{eq:5a} fails the spectral
question is not asked.

\emph{Complexity.} Every step is a rank computation or an exact linear solve on a
matrix of at most $2\omega+1$ rows and $q$ columns whose entries have bit length
$\operatorname{poly}(n,\mathcal L)$ by Lemma~\ref{lem:coeff}, so Gaussian
elimination over $\mathbb Q$ on such a matrix costs $\operatorname{poly}(n,\mathcal L)$
bit operations; step 4 costs
$\operatorname{poly}(q,\omega,\mathcal L_0)\subseteq\operatorname{poly}(n,\mathcal L)$
by Lemma~\ref{lem:zero}.
\end{proof}

\begin{remark}[Kronecker canonical form]\label{rem:kronecker}
The pencil $Q$ has full column rank at every finite complex value if and only if its Kronecker
form (in the sense of \cite{bee88}) contains no right singular block and no finite regular
elementary-divisor block. Left singular blocks and \emph{infinite} regular elementary divisors are
permitted, since they produce no rank drop at a finite $\lambda$: for instance
$Q(\lambda)=1$, $Q(\lambda)=I_2$ and $Q(\lambda)=\begin{bmatrix}1&\lambda\\0&1\end{bmatrix}$ all
satisfy the condition. Only the \emph{finite} spectrum is constrained.

The reason is that the condition is a statement about the rank function on the finite line: a block
contributes to a violation exactly when it makes the rank fall below the number of columns at some
finite value, which is what a right singular block does at every finite value and what a finite
elementary divisor does at its own root; a left singular block, by contrast, is of full row rank, so it
adds rows without adding deficiency, and an infinite elementary divisor acts at $\lambda=\infty$, a
point that \eqref{eq:5b} does not evaluate. This is also why the pencil of \eqref{eq:5b} is allowed to
carry structure that a square nonsingular pencil could not: with $q>\omega$ the factor $H$ has a
constant left null vector, and every such vector is a left null vector of $\lambda X-G$ as well.
\end{remark}

\begin{remark}[Witness]\label{rem:witness}
If \eqref{eq:R0} is consistent, exact elimination returns
$S_0,\dots,S_{2\omega-2}$, hence an explicit polynomial left inverse
$S(\lambda)$, verifiable by multiplying the two polynomial matrices. The test is
therefore not merely a yes/no decision.
\end{remark}

\section{Membership in the existential theory of the reals}\label{sec:existsr}

This section places the decision problem at a prescribed order in the existential theory of the reals, and records the time bound that follows from the two algorithms.

\begin{thm}\label{thm:existsr}
For every fixed $q$, $\mathrm{DEC}(q)\in\exists\mathbb R$. Consequently the threshold problem $\mathrm{MINDEC}(Q)$ is in $\exists\mathbb R$ for every $Q$.
\end{thm}

\begin{proof}
Appendix~\ref{app:encoding} constructs, for each $\omega=0,\dots,\min(q,k)$, an existential formula
$\Phi_\omega$ whose size is counted there: $O(n^3)$ real variables and $O(n^3)$ equations of degree at
most two, with no enumeration of column sets. The disjunction
\begin{equation}\label{eq:existsr-disj}
\Phi=\bigvee_{\omega=0}^{\min(q,k)}\Phi_\omega
\end{equation}
has $O(k)$ branches of polynomial size, is built in polynomial time from the input, and its solutions correspond exactly to feasible $L$:
\begin{equation*}
\begin{aligned}
&\Phi\ \text{satisfiable}\\
&\iff\exists\,L=[L_0;R]\ \text{satisfying}\ \eqref{eq:5a}\ \text{and}\ \eqref{eq:5b}.
\end{aligned}
\end{equation*}
Thus $\Phi$ is a polynomial-time reduction to the existential theory of the reals, so
$\mathrm{DEC}(q)\in\exists\mathbb R$; membership is a statement about the reduction, and decidability of
the target theory is not used. For $\mathrm{MINDEC}(Q)$ take the finite disjunction \eqref{eq:mindec}
over $q=r_0,\dots,\min(Q,n)$, which is a reduction of the same kind; hence $\mathrm{MINDEC}(Q)\in
\exists\mathbb R$ for every $Q$.
\end{proof}

\begin{cor}\label{cor:time}
For every $r\ge0$, the decision time in the bit model satisfies
\begin{equation*}
T\le\min\bigl\{\mathcal L^{O(q^2)},\
\exp\bigl(O(\ell^2\log(\ell+2))\bigr)\operatorname{poly}(\mathcal L)\bigr\}.
\end{equation*}
\end{cor}

\begin{proof}
By \eqref{eq:fpt-l} and the bound of Theorem~\ref{thm:fixedq}, and since both algorithms decide the
same language, the smaller bound applies.
The exponent in the second bound does not shrink when $r=0$, since $\ell$ does not involve $r$; what
happens for $r=0$ is that Section~\ref{sec:zero} supplies a third algorithm, of cost
$\operatorname{poly}(\mathcal L)$, so the minimum in the statement is polynomial there.
\end{proof}


\section{Conclusions}\label{sec:conclusions}

The two upper bounds are parameterized, and the parameter enters a different part of each: $q$
governs the exponent of $\mathcal L^{O(q^2)}$, while $\ell=\operatorname{rank}(AN)$, which does not
involve the number $r$ of free rows, governs the factor
$\exp\bigl(O(\ell^2\log(\ell+2))\bigr)$ of Corollary~\ref{cor:time}. With no free rows the bound is
$\operatorname{poly}(\mathcal L)$ (Section~\ref{sec:zero}). Membership in the existential theory of
the reals holds for the prescribed-order problem and for its finite threshold version,
$\mathrm{DEC}(q),\mathrm{MINDEC}(Q)\in\exists\mathbb R$ (Theorem~\ref{thm:existsr}). The same
reduction underlies the minimum-order and construction results, which are treated in the companion
paper.
The feasibility characterization itself is taken from \cite{fer10,dar25,fer26a,fer26b}; what is new here is the compression to $O(q^2)$ variables, independently of the state dimension, and the complexity bounds that follow from it.

\subsection{Limitations and what is not claimed}\label{sec:limits}

We do not claim that $\mathrm{DEC}\in\mathrm P$, that it is $\exists\mathbb R$-complete, or any
hardness result; hardness and gadget constructions are outside our scope. Neither parameter is
bounded in the unrestricted problem, so the two bounds are compatible with all three outcomes ---
$\mathrm{DEC}\in\mathrm P$, $\exists\mathbb R$-completeness, or an intermediate parameterized
classification --- and the paper claims none of them: the unrestricted complexity of $\mathrm{DEC}$
remains open. The alternatives concern $\mathrm{DEC}$; a polynomial-time algorithm for it would
also settle the minimum-order question treated in the companion paper.

Two items remain outside the verified core.
\begin{enumerate}
\item \textbf{The $r=0$ branch: resolved, with a residual audit.} For $r=0$ the two conditions reduce to a constant-rank test and to the constancy of the rank of the degree-one pencil $\lambda X-G$. Section~\ref{sec:zero} decides the latter by an exact, deterministic rational linear system of
polynomial size, read off the degree bound $2\omega-2$; it returns an explicit witness. What remains
is the line-by-line audit of the Hadamard and elimination constants in its bit bound.
\item \textbf{Coefficient bounds.} The coefficient-size lemma in Appendix~\ref{app:bits} is paper-level; its constants have not been independently audited.
\end{enumerate}

\subsection{The size of the parameter}\label{sec:paramsize}

Two quantities describe the free data: the number of coordinates
\begin{equation}\label{eq:kappa-size}
\kappa_{\rm obs}=r\,d_0\le2rn
\end{equation}
needed to write the $r$ free rows, and the parameter that governs the bound,
\begin{equation*}
\ell=\operatorname{rank}(AN)\le k\le n .
\end{equation*}
Only the first can be quadratic in $n$. By the subspace reduction of Section~\ref{sec:subspace}
(Lemmas~\ref{lem:struct} and~\ref{lem:sandwich}) the bound of Theorem~\ref{thm:fpt} depends on
$\ell$ and not on $\kappa_{\rm obs}$: the coordinate count \eqref{eq:kappa-size} records how the
free rows are written, not how hard the instance is.

Whether $\ell$ is optimal is open. Two further reductions are conceivable, and they are of
different kinds.

\begin{enumerate}
\item \emph{Quotients of the observer rows.} The conditions depend on row spaces, hence are
invariant under row operations of $L$. Those preserving $L_0$ are block lower-triangular,
\begin{equation*}
T=\begin{bmatrix}I_{r_0}&0\\ U&S\end{bmatrix},\qquad \det S\ne0 ,
\end{equation*}
so two different free blocks can present the same $W$. The reduction of
Section~\ref{sec:subspace} already quotients by this symmetry; the question is whether a
smaller invariant than $\ell$ exists.
\item \emph{Spectral invisibility.} One might hope to divide by directions that change neither
the rank condition nor the finite spectrum. Such a quotient must preserve the interaction
between $H$ and $H_A$, not only $\operatorname{rank}H$; a quotient preserving \eqref{eq:5a} but
changing the spectrum would be invalid. We do not know whether one exists.
\end{enumerate}

\paragraph*{Growth of the parameter-dependent factor}
By \eqref{eq:nvar}, $sD=O(\ell^2)$ and $N_{\rm var}\le\ell^2+2\ell$, so
$f(\ell)=\exp\bigl(O(\ell^2\log(\ell+2))\bigr)$. The result is primarily structural: it
identifies an invariant of the instance, independent of how many free rows are allowed, that
governs the complexity bound, although the doubly exponential dependence on $\ell$ may limit
practical applicability for large values.

\subsection{Reproducibility and evidence}\label{sec:repro}

The results rest on three kinds of evidence, and the accompanying artifact records which statement rests on which.

\begin{enumerate}
\item \emph{Lean~4 with Mathlib.} The labels of Table~\ref{tab:evidence} separate statements
formalized as stated (\emph{literal}) from those formalized at the level of linear maps and
submodules (\emph{abstract core}), the translation to matrices being paper-level; the single
\emph{ingredient} is one step of the proof of the certificate. Every declaration is checked by the
kernel against only (\texttt{propext}, \texttt{Classical.choice}, \texttt{Quot.sound}). The
reductions, the certificate and the complexity accounting are paper-level, so this is not a formal
verification of the algorithmic theorems.
\item \emph{Paper-level proofs.} The Cauchy--Binet compression of the certificate, the coefficient-size estimates and the existential-real encoding.
\item \emph{Exact rational arithmetic.} The identities on which the constructions rest are re-checked, both on randomly generated instances and on the constructed families; this excludes arithmetic slips but does not replace the proofs.
\end{enumerate}

The Lean development and the exact-arithmetic scripts form the artifact at
\url{https://github.com/Zhengshuji/q1-functional-observer}. The matrix-level identifications the Lean
core does not cover --- the pencil/gcd bridge, the matrix form of $\mathcal T$, the converse step from
a polynomial left inverse to \eqref{eq:R0}, and the encoding --- are paper-level and are the ones
re-checked by exact arithmetic. Of the classical facts invoked by citation, the rank identity, the
solvability criterion and the B\'ezout steps of Lemmas~\ref{lem:rankid}, \ref{lem:rankcons}
and~\ref{lem:bezout} are also proved in the artifact, so there the citation only keeps the proofs
short; the remaining cited facts are not formalized.

\begin{table*}[!t]
\centering
\caption{Where each statement of the paper is verified. \emph{Literal}: the Lean statement is the
statement of the paper. \emph{Abstract core}: the Lean statement is the linear-map or submodule
version, and the translation to matrices is paper-level. \emph{Ingredient}: one step of the proof,
not the statement. Statements not listed are paper-level throughout.}
\label{tab:evidence}
\footnotesize
\begin{tabular}{@{}lll@{}}
\toprule
Statement & Lean evidence & Formal name \\
\midrule
Rank identity (Lemma~\ref{lem:rankid}), \eqref{eq:rankid} & literal &
\texttt{rank\_vstack\_add\_rank\_mul} \\
Compression criterion (Lemma~\ref{lem:compression}) & abstract core &
\texttt{ker\_compressed\_iff} \\
Space of row pairs (Lemma~\ref{lem:attainable}) & abstract core & \texttt{freePairSpace},
\texttt{kappaObs\_eq\_mul\_freePairSpace} \\
Sandwich and realizability (Lemma~\ref{lem:sandwich}) & literal &
\texttt{rowSpan\_fromRows\_left}, \texttt{mem\_rowSpan\_iff\_exists\_vecMul} \\
Conditions depend on $W$ (Lemma~\ref{lem:W-conditions}) & abstract core &
\texttt{rank\_fromRows\_eq\_rank\_iff}, \texttt{rowSpace\_eq} \\
Structural bound, upper half (Lemma~\ref{lem:struct}) & literal &
\texttt{finrank\_le\_two\_mul\_of\_pi1\_eq\_pi2} \\
Structural bound, lower half (Lemma~\ref{lem:struct}) & literal &
\texttt{finrank\_add\_one\_le\_of\_spectral} \\
Rank consistency (Lemma~\ref{lem:rankcons}) & literal &
\texttt{exists\_mulVec\_eq\_iff\_rank\_gram} \\
B\'ezout with degree bound (Lemma~\ref{lem:bezout}) & literal & \texttt{exists\_bezout\_family} \\
Rectangular left inverse (Proposition~\ref{prop:zero}) & literal &
\texttt{exists\_leftInverse\_linear} \\
Gram certificate, kernel identity (Proposition~\ref{prop:gram}) & ingredient &
\texttt{sum\_sq\_zero\_iff} \\
Coefficient system \eqref{eq:R0} $\Rightarrow$ left inverse & literal (one direction) &
\texttt{spoly\_mul\_pencil\_eq\_one} \\
Matrix-product coefficients & literal & \texttt{coeff\_mul\_apply} \\
Degree bounds $\deg\det$, $\deg\operatorname{adj}$ & literal &
\texttt{natDegree\_det\_le}, \texttt{natDegree\_adjugate\_le} \\
Zero-free-row test (Section~\ref{sec:zero}) & paper-level & --- \\
Encoding (Theorem~\ref{thm:existsr}) & paper-level & --- \\
\bottomrule
\end{tabular}
\end{table*}

\appendices
\section{Bit sizes and the coefficient bound}\label{app:bits}

We use the standard bit model: the input has rational entries of total bit length $\mathcal L$,
and $n,p,q\le\mathcal L$. Since only the row space of $C$ enters both conditions, we replace $C$ by
a rational basis of its row space when it has more than $n$ rows; this preserves $\operatorname{row}C$
and $\operatorname{row}[C;CA]$, hence $K_0$ and \eqref{eq:5a}--\eqref{eq:5b}, costs
$\operatorname{poly}(\mathcal L)$ bit operations, and leaves an output matrix with at most $n$ rows.
We therefore assume $p\le n$ in what follows. Renegar's Part II bit bound
\cite{ren92} has the form
\begin{equation}\label{eq:rengar}
\beta(\log\beta)(\log\log\beta)(sD)^{O(N_{\rm var})},
\end{equation}
for integer-coefficient input, where $N_{\rm var},s,D,\beta$ are the number of real variables, the number of
polynomials, the maximum degree and the maximum coefficient bit length. Only $(sD)^{O(N_{\rm var})}$ is
exponential; the dependence on $\beta$ is polynomial in $\log\beta$. We apply \eqref{eq:rengar} in the
following normalization, used throughout: every system is cleared of denominators (the solution set
and the degree are unchanged and $\beta$ is inflated by the common denominator only); equations and
inequalities count once each in $s$; a strict $g>0$ is replaced by $g\ge\varepsilon$ with one fresh
variable $\varepsilon>0$, adding at most $s$ variables and leaving $s,D$ fixed; and the bound covers
the decision only, the construction being accounted for by Lemmas~\ref{lem:rationalbasis}
and~\ref{lem:coeff}.

The estimates below are the ones the bound \eqref{eq:fpt-l} and the counting of
Section~\ref{sec:fixedq} rely on. Each is stated with the mechanism that produces it and with its audit
status: an entry marked \emph{Hadamard} rests on the Hadamard bound for a determinant of a matrix whose
entry bit length is already known, so the dependence on $\mathcal L$ is polynomial but the elementary
constant has not been re-derived here.

\begin{table}[!t]
\centering
\caption{Bit lengths of the objects entering the bound. $n,p,q\le\mathcal L$ and every entry of the
input has bit length at most $\mathcal L$; $\tau$ denotes the entry bit length of the matrix at the
stage concerned.}
\label{tab:bitlen}
\footnotesize
\begin{tabular}{@{}p{0.30\columnwidth}p{0.30\columnwidth}p{0.26\columnwidth}@{}}
\toprule
Object & Bit length & Mechanism \\
\midrule
$C$, $A$, $L_0$ & $\le\mathcal L$ & input \\
kernel basis $N$ of $[C;CA]$ & $O\bigl(k(\tau+\log n)\bigr)\le\operatorname{poly}(\mathcal L)$ &
Hadamard + Cramer \\
$M=[N\ AN]$ & $O(\tau_M+\log n)$ & submultiplicativity \\
$D$ with $\ker D=\operatorname{row}M$ & $O\bigl(d_0(\tau_M+\log k)\bigr)$ & Hadamard + Cramer \\
pivot block $X$ of $H=XB$ & $\operatorname{poly}(\mathcal L)$ & fraction-free elimination,
entry growth $O(n)$ per step \\
windows $Z$ of the minors of $\lambda X-G$ & $\operatorname{poly}(\mathcal L)$ & sum over
$\binom{2\omega}{\omega}$ products of $\omega$ entries \\
$Z^{\mathsf T}Z$ & $\operatorname{poly}(\mathcal L)$ & Cauchy--Binet, $\binom{2\omega}{\omega}$
terms \\
$e_\omega(X^{\mathsf T}X)$, $e_0$ & $O(\omega\log\omega)$ above $\tau_Z$ & binomial coefficients,
$\lambda$-free \\
cleared system $\Psi$ & $\beta=\operatorname{poly}(\mathcal L)$ & one common denominator per
equation \\
\bottomrule
\end{tabular}
\end{table}

The two entries that need care are the pivot block and the Cauchy--Binet compression. The first comes
from Gaussian elimination on a rational matrix of size at most $n\times k$ whose entries are rational
functions of the input: with fraction-free elimination the intermediate numerators grow by at most one
factor of a subdeterminant per step, so the growth is $O(n)$ in the exponent and the entries stay of
poly$(n,\mathcal L)$ bit length \cite[Ch.~6]{vzgg13}. The second is a sum of $\binom{2\omega}{\omega}$
products of $\omega$ minors, and each minor is itself a determinant of a $2\omega\times2\omega$ matrix
with entries of bit length $\tau_Z$; Hadamard gives $\omega\tau_Z+O(\omega\log\omega)$ for the minor and
the sum adds $\log\binom{2\omega}{\omega}$ bits, which is $O(\omega)$ by Stirling. Both are polynomial
in $\mathcal L$ because $\omega\le k\le n\le\mathcal L$.

\begin{lem}[Rational bases and annihilators]\label{lem:rationalbasis}
Let $M\in\mathbb Q^{m\times n}$ have entries of bit length at most $\tau$, let $r=\operatorname{rank}M$
and put $\tau_M=2r(\tau+\log(r+1)+1)+1$. Then a basis $N$ of $\ker M$, a basis $R$ of
$\operatorname{row}M$, and a matrix $D\in\mathbb Q^{(n-r)\times n}$ with
$\ker D=\operatorname{row}M$ can be computed in $\operatorname{poly}(m,n,\tau)$ bit operations, and
the entries of $N$, $R$ and $D$ have bit length at most $\tau_M$.
\end{lem}

\begin{proof}
Run fraction-free (Bareiss) Gaussian elimination on $M$ and keep the pivot rows and columns. Every
entry of the resulting reduced row echelon form is a ratio of minors of $M$: by Cramer's rule the
determinant in the numerator is an $s\times s$ minor with $s\le r$, and by Hadamard's inequality,
for $s\ge1$,
\begin{equation*}
\operatorname{bits}\det M_{I,J}
\le s\bigl(\tau+\tfrac12\log s\bigr)+1
\le r\bigl(\tau+\tfrac12\log r\bigr)+1\le\tau_M .
\end{equation*}
Fraction-free elimination carries a single minor as common denominator at every stage, of size at
most $r\times r$; the entries of $N$, $R$ and $D$ are therefore ratios of two $r\times r$ minors of
$M$, of bit length at most $\tau_M$. The pivot columns give $R$ and the free columns give $N$ after
back-substitution; put $D=N^{\mathsf T}$. The rows of $D$ are the columns of $N$, hence a basis of
the right null space $\{y\in\mathbb Q^n:My=0\}$ of $M$, which is the orthogonal complement of
$\operatorname{row}M$; thus $\ker D=\operatorname{row}M$, and the bit length of $D$ is that of $N$.
The elimination uses
$O(mnr)$ arithmetic operations on numbers of the displayed bit length, so the total is
$\operatorname{poly}(m,n,\tau)$ bit operations.
\end{proof}

\begin{lem}[Coefficient sizes]\label{lem:coeff}
Let the input entries have bit length at most $\mathcal L$, and let
$n\le\mathcal L$. Then:
\emph{(i)} the kernel basis $N$, the matrix $D$, the pivot block $X$ and the block $G$ of the
rank factorizations produced by the elimination of Section~\ref{sec:pivot}, the representative pair
$H=XB$, $H_A=GB$ of such a factorization, and the prescribed products $h^0=L_0N$, $g^0=L_0AN$ have
entries of bit length $\operatorname{poly}(n,\mathcal L)$;
\emph{(ii)} every maximal minor $f_i$ of $P(\lambda)=\lambda H-H_A$ of the representative pair in
(i), and every coefficient of $f_i$, has bit length $\operatorname{poly}(n,\mathcal L)$;
\emph{(iii)} the entries of $Z^{\mathsf T}Z$ have bit length
$\operatorname{poly}(n,\mathcal L)$.
\end{lem}

\begin{proof}
Each object bounded below is a ratio of minors of the integer input, or a short sum of products of two
such minors, so the Hadamard estimate of Lemma~\ref{lem:rationalbasis} bounds its entries.

\emph{(i)} Apply Lemma~\ref{lem:rationalbasis} to $[C;CA]$, whose entries have bit length at most
$\mathcal L$ and whose rank is $n-k$: the basis $N$ of $\ker[C;CA]$ has entries of bit length
$O\bigl((n-k)(\mathcal L+\log n)\bigr)$. Applying the lemma again to $[N\ AN]$, whose entries now
have that bit length and whose rank is $d_0$, gives the matrix
$D\in\mathbb Q^{(2k-d_0)\times2k}$ with $\ker D=\operatorname{row}[N\ AN]$, with entries of
bit length $O\bigl(d_0((n-k)(\mathcal L+\log n)+\log(k+1))\bigr)=\operatorname{poly}(n,\mathcal L)$;
the two eliminations cost $\operatorname{poly}(\mathcal L)$ bit operations.
The pivot block $X$, the block $G$ and the coordinate matrix $B$ are produced by the same
elimination whenever the compressed pair is prescribed, in particular for $r=0$, where $H=L_0N$ and
$H_A=L_0AN$: then $X$ is the pivot block of $H$, $B$ is the coefficient matrix with
$B_{:,J}=I_\omega$, and $G=FX$ for the rational solution $F$ of $FH=H_A$, so all three are
submatrices or ratios of minors of input-derived matrices. For $r\ge1$ the blocks
$X,G$ are the quantified variables of \eqref{eq:phi} and enter its constraints as indeterminates. No bound on
the entries of $H=LN$ for an arbitrary feasible $L$ is asserted. The
prescribed products are sums of $n$ products of an input entry with an entry of $N$, hence have bit
length $\mathcal L+\operatorname{poly}(n,\mathcal L)+O(\log n)=\operatorname{poly}(n,\mathcal L)$.

\emph{(ii)} A maximal minor $f_i$ of $P(\lambda)=\lambda H-H_A$ is the determinant
of an $\omega\times\omega$ matrix whose entries are linear in those of $H,H_A$.
Laplace expansion along the $a$ columns that carry the factor $\lambda$ shows that the coefficient
of $\lambda^a$ is a signed sum over the $\binom{\omega}{a}$ index sets $S\subseteq[\omega]$ with
$|S|=a$ of $\det H_{I,S}\det H_{A,I,S^{c}}$, and Hadamard's inequality bounds each of the two
factors, hence each such product, by $O\bigl(\omega(\operatorname{bits}H+\log\omega)\bigr)$ bits:
\begin{equation*}
\deg f_i\le\omega,\qquad
\operatorname{bits}c_{i,a}=\operatorname{poly}(n,\mathcal L).
\end{equation*}

\emph{(iii)} The case $\omega=0$ gives $Z^{\mathsf T}Z=0$; assume $\omega\ge1$. By
item~\ref{item:coef} of Section~\ref{sec:certificate},
$(Z^{\mathsf T}Z)_{t,t'}=\sum_{j=0}^{\omega-1}S_{t-j,t'-j}$, where each
$S_{a,b}=\sum_ic_{i,a}c_{i,b}$ is a sum of $\#\{I\}=\binom q\omega\binom k\omega$ products of two
coefficients. Hence $(Z^{\mathsf T}Z)_{t,t'}$ is a sum of at most $\omega\cdot\#\{I\}$ terms, each
of bit length at most $2\operatorname{poly}(n,\mathcal L)$. A sum of $T$ terms of bit length $b$ has
bit length $b+\log T$, and
\begin{equation*}
\log\bigl(\omega\cdot\#\{I\}\bigr)\le\omega+2\omega\log n=O(n),
\end{equation*}
so
\begin{equation*}
\operatorname{bits}(Z^{\mathsf T}Z)_{ij}
\le2\operatorname{poly}(n,\mathcal L)+O(n)
=\operatorname{poly}(n,\mathcal L).
\end{equation*}
\end{proof}

It remains to say why the per-chart bound already gives \eqref{eq:fpt-l}: the enumeration of
$t$, of the charts and of $\omega$ contributes the factors $2\ell+1$, $4^{\ell}$ and $\ell$, all
absorbed by the exponential; every chart and every $\omega$ costs one call to the decision procedure
of \eqref{eq:nvar}, and the rational data of each call have bit length
$\operatorname{poly}(n,\mathcal L)$ by Lemma~\ref{lem:coeff}. The cases $k=0$ and $\omega=0$
are treated directly in Section~\ref{sec:deg}, and $r=0$ in Section~\ref{sec:zero}. For fixed $q$
the sharper parameter is $q$ itself, and the bound $\mathcal L^{O(q^2)}$ of
Section~\ref{sec:fixedq} applies instead.

\section{The existential-real encoding}\label{app:encoding}

This appendix constructs the encoding used in the proof of Theorem~\ref{thm:existsr}. Throughout, $q$ is fixed and the notation is that of Section~\ref{sec:compressed}: $N$ is a basis matrix of $\ker[C;CA]$, $H=LN$, $H_A=LAN$, and $\omega=\operatorname{rank}H$.

For fixed $q$, \eqref{eq:5b} is the condition that the $q\times\omega$ pencil
\begin{equation*}
Q(\lambda)=\lambda X-G
\end{equation*}
has rank $\omega$ for every $\lambda\in\mathbb C$. This is stronger than having rank $\omega$ over the fraction field $\mathbb C(\lambda)$: the latter does not give $\operatorname{rank}Q(\lambda)=\omega$ at every $\lambda$, as
\begin{equation*}
Q(\lambda)=\begin{bmatrix}\lambda\\ \lambda\end{bmatrix},\qquad
\operatorname{rank}_{\mathbb C(\lambda)}Q=1,\qquad \operatorname{rank}Q(0)=0
\end{equation*}
shows.

This is the pencil of Lemma~\ref{lem:zero}: by its parts (i)--(iv) and the Quillen--Suslin/B\'ezout
criterion over a principal ideal domain \cite{fab90,lan02}, the condition is equivalent to a
polynomial left inverse $S$ with $\deg S\le2\omega-2$, which the encoding below never forms.

For the $\exists\mathbb R$ encoding the minors are not enumerated. For each $\omega=0,\dots,\min(q,k)$ introduce the blocks
\begin{equation*}
\begin{gathered}
X\in\mathbb R^{q\times\omega},\qquad B\in\mathbb R^{\omega\times k},\\
G\in\mathbb R^{q\times\omega},\qquad X_L\in\mathbb R^{\omega\times q},\\
B_R\in\mathbb R^{k\times\omega},
\end{gathered}
\end{equation*}
and the coefficients of $S(\lambda)=\sum_{j=0}^{2\omega-2}S_j\lambda^j$, where $S_j\in\mathbb R^{\omega\times q}$, and impose
\begin{equation*}
\begin{gathered}
H=XB,\qquad H_A=GB,\\
X_LX=I_\omega,\qquad BB_R=I_\omega,\\
S(\lambda)(\lambda X-G)=I_\omega .
\end{gathered}
\end{equation*}
The polynomial identity is linear in the coefficients $S_j$; comparing coefficients of $\lambda^j$ makes it the linear system \eqref{eq:R0}. The rank conditions $X_LX=I_\omega$ and $BB_R=I_\omega$ are polynomial equations, and $BB_R=I_\omega$ makes $B$ full row rank, hence, as in Section~\ref{sec:spectral}, $\operatorname{rank}P(\lambda)=\operatorname{rank}(\lambda X-G)=\omega$ for all $\lambda$. Also $H_A=GB=(GX_L)XB=FH$ with $F=GX_L$, so \eqref{eq:5a} holds.

It remains to tie $(H,H_A)$ to the input. The rows of $H=LN$, $H_A=LAN$ satisfy, with $h_i^0=(L_0N)_i$ and $g_i^0=(L_0AN)_i$,
\begin{equation}
X_iB=h_i^0,\qquad G_iB=g_i^0\qquad(1\le i\le r_0),
\tag{A.1}\label{eq:B1}
\end{equation}
and, for the free rows, $(H_i,H_{A,i})\in\mathcal T=\operatorname{row}[N\ AN]$, equivalently
\begin{equation}
D\begin{bmatrix}(X_iB)^{\mathsf T}\\ (G_iB)^{\mathsf T}\end{bmatrix}=0
\qquad(r_0<i\le q),
\tag{A.2}\label{eq:B2}
\end{equation}
where $D$ is a rational matrix with $\ker D=\mathcal T$. Both are polynomial (in fact bilinear) equations. They are not optional: without \eqref{eq:B1}--\eqref{eq:B2} the encoding would accept arbitrary pairs $(H,H_A)$.

The correctness splits into two independent statements.

\paragraph*{Soundness}
Suppose the system has a solution. Set $L$ with first $r_0$ rows $L_0$ and, for each $i>r_0$, a preimage $r_i$ with $(X_iB,G_iB)=r_i[N\ AN]$, which exists by \eqref{eq:B2}. Then $H=LN$, $H_A=LAN$ and $H_A=GB=(GX_L)XB=FH$ with $F=GX_L$, so \eqref{eq:5a} holds; and $\operatorname{rank}(\lambda H-H_A)=\operatorname{rank}(\lambda X-G)=\omega$ for all $\lambda$, which gives \eqref{eq:5b} and uses $BB_R=I_\omega$ (full row rank of $B$) together with $S(\lambda)(\lambda X-G)=I_\omega$.

\paragraph*{Completeness}
A feasible $L$ supplies $X,B,X_L,B_R$ from a rank factorization $H=XB$ with $B_{:,J}=I_\omega$, $G=FX$ with $H_A=FH$ (Lemma~\ref{lem:compression}), and $S$ from the left-inverse construction above; then \eqref{eq:B1} and \eqref{eq:B2} hold by construction, and the remaining equations are identities.

The $\omega=0$ case is the branch $H=H_A=0$, which \eqref{eq:B1} forces when $L_0N=0$ and rejects otherwise; it is a separate algebraic branch, not a degree-negative instance of \eqref{eq:R0}. The number of variables is
\begin{equation*}
O(q\omega+\omega k+q\omega+\omega q+k\omega+(2\omega-1)\omega q)
=O(n^3),
\end{equation*}
the number of equations is polynomial, and the degree of the left inverse is $2\omega-2\le2n$. The disjunction \eqref{eq:existsr-disj} over $\omega$ has at most $n+1$ branches. Hence the encoding has polynomial size. This proves $\mathrm{DEC}(q)\in\exists\mathbb R$ for every fixed $q$, and hence $\mathrm{MINDEC}(Q)\in\exists\mathbb R$.

\begin{thebibliography}{99}
\bibitem{fer10} J. A. Fernando, H. Trinh, L. Jennings. Functional observability and the design of minimum order linear functional observers. \emph{IEEE Trans. Automat. Control} 55(5):1268--1273, 2010.
\bibitem{dar25} M. Darouach, J. A. Fernando. On functional observability and functional observer design. \emph{Automatica}, vol. 173, Art. no. 112115, 2025, doi: 10.1016/j.automatica.2025.112115.
\bibitem{fer26a} J. A. Fernando. A minimum-order functional observer beyond the Darouach and Luenberger constructions. arXiv:2609.33615, 2026 (preprint).
\bibitem{fer26b} J. A. Fernando. Complete characterization of minimum-order functional observers from Darouach to Luenberger. arXiv:2609.37992, 2026 (preprint).
\bibitem{zha25} Y. Zhang, T. Fernando, M. Darouach. Functional
  observability, structural functional observability, and optimal sensor
  placement. \emph{IEEE Trans. Automat. Control}, vol. 70, no. 3,
  pp. 1592--1607, 2025, doi: 10.1109/TAC.2024.3462558.
  (arXiv:2307.08923.)
\bibitem{dey21} P. Dey, N. Balachandran, D. Chatterjee. Complexity of
  constrained sensor placement problems for optimal observability.
  \emph{Automatica}, vol. 131, Art. no. 109758, 2021,
  doi: 10.1016/j.automatica.2021.109758.
\bibitem{lof26} J. L\"{o}fberg. Static output-feedback stabilization is
  NP-hard. arXiv:2609.16886, 2026 (preprint).
\bibitem{bar26} G. Barkai, I. Shames. On the NP-hardness of unconstrained
  static output feedback stabilization. arXiv:2609.38728, 2026 (preprint).
\bibitem{liu26} Z. Liu. Computational complexity of static output feedback
  stabilization. arXiv:2610.00390, 2026 (preprint).
\bibitem{maj26} A. Majumdar. Affine rank minimization is
  $\exists\mathbb R$ complete. arXiv:2602.14037, 2026 (preprint).
\bibitem{hj13} R. A. Horn, C. R. Johnson. \emph{Matrix Analysis}. 2nd ed.
Cambridge University Press, 2013.
\bibitem{vzgg13} J. von zur Gathen, J. Gerhard. \emph{Modern Computer Algebra}.
3rd ed. Cambridge University Press, 2013.
\bibitem{emi99} I. Z. Emiris, B. Mourrain. Matrices in elimination theory. \emph{J. Symbolic Comput.} 28(1--2):3--44, 1999, doi: 10.1006/jsco.1998.0266.
\bibitem{lan02} S. Lang. \emph{Algebra}. Revised 3rd ed.
Graduate Texts in Mathematics 211, Springer, 2002.
\bibitem{fab90} A. Fabiano, D. C. Struppa. Bounds for the degrees of the entries of left
inverses of polynomial matrices. \emph{Multidimensional Systems and Signal Processing},
vol.~1, no.~1, pp.~67--74, 1990, doi: 10.1007/BF01812207.
\bibitem{bee88} Th. Beelen, P. Van Dooren. An improved algorithm for the computation of Kronecker's canonical form of a singular pencil. \emph{Linear Algebra Appl.} 105:9--65, 1988.
\bibitem{ren92} J. Renegar. On the computational complexity and geometry of the first-order theory of the reals. Part II: The general decision problem. \emph{J. Symbolic Comput.} 13(3):301--327, 1992, doi: 10.1016/S0747-7171(10)80004-5.
\bibitem{companion2} Minimum order and minimum-order construction for functional observers.
Companion manuscript, 2026. Cited for the optimization and construction problems and for the
spectral lower bound on the minimum order.
\end{thebibliography}

\end{document}
